% Commentary on manuscript pp. 138--161. Anchors refer to the supplied transcription.

\begin{dcomment}{exposition}{p138:Mais on ne va pas faire bouger}{The base changes before the representative does}
Grothendieck has reached the point at which the categories constructed earlier must begin to behave as a theory of coefficients. A category over $S$ is meant to represent a homotopy type varying over $S$; moving it to another base should therefore support the same adjunction as restriction and left Kan extension of diagrams. The difficulty is concentrated in restriction. Composition of structure maps gives $f_!$ directly, whereas a categorical pullback computes $f^*$ only after a suitable replacement. Keeping these two facts together explains the apparently elaborate calculation that follows.

Write $Q_S:\Cat/S\longrightarrow\HOT_W(S)$ for the localization functor, and let $p:X\longrightarrow S$ be a representative. The replacement $\Phi_S(X)$ has objects
\[
 (s,x,a),\qquad a:s\longrightarrow p(x),
\]
and a morphism $(s,x,a)\longrightarrow(t,y,b)$ consists of $u:s\longrightarrow t$ and $v:x\longrightarrow y$ with $p(v)a=bu$. Its projection to $S$ is a split fibration: a chosen lift of $u:s\longrightarrow t$ with target $(t,y,b)$ is $(s,y,bu)\longrightarrow(t,y,b)$. The fibre at $s$ is the comma category $s\backslash X=(s\downarrow p)$. The natural functor
\[
 i_X:X\longrightarrow\Phi_S(X),\qquad
 x\longmapsto(p(x),x,\id_{p(x)}),
\]
is a colocally $W$-equivalent replacement, as established in the preceding construction. Thus the usable formulas are
\[
 f_!Q_{S'}(X')=Q_S(X'),\qquad
 f^*Q_S(X)=Q_{S'}(\Phi_S(X)\times_S S').
\]
The first $X'$ has merely changed its structure map. In the second formula the replacement is made over the base $S$ from which we are restricting.

This order is essential. Even if $X'\longrightarrow S'$ is a fibration, the composite $X'\longrightarrow S$ need not be one. For instance, the identity of the point $\{1\}$ is a fibration, but its composite with $\{1\}\hookrightarrow[1]=(0\to1)$ is not: the arrow $0\to1$ has no lift with target the sole object. Replacing over $S'$ and then composing with $f$ therefore does not solve the problem. On the other hand, a smooth representative $X\longrightarrow S$ may already be pulled back, and the comparison with the formula using $\Phi_S(X)$ is an isomorphism after localization.

There is also a useful variance check. In the usual homotopy theory $W=W_\infty$, a fibration over $S$ describes a contravariant diagram of spaces. The transition along $s\to t$ runs from the fibre at $t$ to the fibre at $s$. The comma category $s\backslash X$, rather than $X/s$, is consequently the value of the associated homotopy diagram at $s$. This is the presheaf convention used in \citet[Introduction, pp.~2--3]{Maltsiniotis2005Asphericite}. It is the opposite of the indexing convention in which a derivator is first presented as $A\mapsto\operatorname{Ho}(\mathcal M^A)$; formulas may be transported between them by replacing every indexing category with its opposite.

The unit is now forced by ordinary pullback. If $r:X'\to S'$ is the structure map, there is a unique functor
\[
 a_{X'}=(i_{X'},r):X'\longrightarrow\Phi_S(X')\times_S S'
\]
whose first component is $i_{X'}$ and whose second is $r$. Naturality of $i$ makes these functors a natural transformation, and the already constructed localized functors turn it into $\alpha:\id\to f^*f_!$. The purpose of the next pages is to construct its partner and verify that the pair is an adjunction; merely recognizing two plausible functors would not supply the necessary compatibility.
\end{dcomment}

\begin{dcomment}{exposition}{p143:ce qui n'est autre que la relation}{A counit that survives localization}
Grothendieck's first attempted counit tries to undo a pullback of the replacement map. The correction on this page is decisive: the inverse must be taken where the map is known to be an equivalence. The natural transformation $i_X$ is invertible in $\HOT_W(S)$, but an arbitrary categorical pullback of $i_X$ need not be invertible. This is a failure of the proposed representative, not a failure of the adjunction.

The rejected construction can be tested with $S=[1]$, $X=\{1\}\to S$, and $S'=\{0\}$. Here $\Phi_S(X)=S$, so the pullback of $i_X$ is $\varnothing\to\{0\}$. For $W_\infty$ this is not a weak equivalence. No inverse to it can be used in defining a counit. This example also explains why a hypothesis on the base change alone will later be a substantive issue.

The corrected counit uses the roof
\[
 \Phi_S(X)\times_S S'
 \xrightarrow{\pi_1}\Phi_S(X)
 \xleftarrow[\ \in W_S^{\mathrm d}\ ]{i_X} X.
\]
In precise notation,
\[
 \beta_{Q_SX}=Q_S(i_X)^{-1}Q_S(\pi_1).
\]
The inverse here exists by construction. Both arrows in the roof vary naturally with $X$, so this formula is natural on representatives and hence on the localization. If $X$ is smooth over $S$, the isomorphism supplied by $i_X\times_S S'$ identifies the roof with the ordinary projection $X\times_S S'\to X$. The two descriptions on manuscript pages 140 and 143 therefore agree; one has not silently changed the counit.

The first triangle identity is now a strict identity followed by cancellation. For $X'$ over $S'$, let $a_{X'}$ be the representative of the unit. By its definition,
\[
 \pi_1a_{X'}=i_{X'}.
\]
After applying $Q_S$ and multiplying by $Q_S(i_{X'})^{-1}$, this becomes $\beta_{f_!Q_{S'}X'}\,f_!(\alpha_{Q_{S'}X'})=\id$. Nothing is required of a pullback of $i_{X'}$ in this argument.

For the second triangle, a smooth representative makes the cancellation legitimate after pullback as well. Take $X$ smooth over $S$, put $X'=X\times_S S'$, and write $r=\pi_1:X'\to X$. The functor
\[
 t=i_X\times_S\id_{S'}:
 X'\longrightarrow\Phi_S(X)\times_S S'
\]
is a colocally $W$-equivalent map over $S'$, because both $X$ and $\Phi_S(X)$ are smooth over $S$. Notice the target: the edition prints $\Phi_S(X')$ once on manuscript page 144, but the map is induced by $i_X$ and must have target $\Phi_S(X)$. The following diagram in the edition already uses the correct target.

Under these identifications the second arrow in the triangle, the image of the counit under $f^*$, is represented in the localization by
\[
 t^{-1}\bigl(\Phi_S(r)\times_S\id_{S'}\bigr):
 \Phi_S(X')\times_S S'\longrightarrow X'.
\]
It is enough to check the identity after composing with $t$. The required equality is
\[
 \bigl(\Phi_S(r)\times_S\id_{S'}\bigr)a_{X'}=t.
\]
Its component in $S'$ says that $a_{X'}$ is a map over $S'$. Its component in $\Phi_S(X)$ is
\[
 \Phi_S(r)i_{X'}=i_Xr,
\]
which is precisely naturality of $i$ at $r$. Both components agree strictly before localization. Since $t$ is invertible afterwards, the second triangle follows.

These computations supply a complete proof, conditional only on the replacement and localization results already in use. They do not require the unproved assertion that all morphisms of $\HOT_W(S)$ are detected as isomorphisms by their fibres. The exact role of the representative is also visible: smoothness is used to pull back one specified weak equivalence in the second triangle, not to make the composite $X'\to S$ smooth. A general formal argument about derived adjunctions packages the same mechanism, but the present calculation records the unit and counit in the form needed for the later base-change maps.
\end{dcomment}

\begin{dcomment}{context}{p146:qui est évidente par fonctorialité}{Integration and the geometry of coefficients}
The adjunction just obtained gives concrete content to the word ``integration'' in Grothendieck's earlier programme. In \emph{Pursuing Stacks}, \S69, he recalls the inadequacy of a single derived category for the trace formalism of SGA~5 and describes the need to retain categories of diagrams together with their operations \citep[\S69, pp.~1--3 of the separately edited text]{Grothendieck1983}. The present construction belongs to the unstable part of that programme: a category over a base carries varying homotopy types, and its total category is the candidate for their homotopical integral.

The analogy becomes an actual universal property here. For the structural functor $p:S\to e$, the operation $p_!$ takes a representative $X\to S$ to the homotopy type of its total category. The adjunction says that maps from this total homotopy type to a coefficient $T$ are the same as maps, over the whole diagram, from the original varying object to the constant object $p^*T$. This is the defining property of a homotopy colimit, expressed at the level of homotopy categories of diagrams. It differs from taking the ordinary colimit of the diagram of already localized fibres: the latter diagram has discarded the coherent data represented by $X$.

The earlier companion's block ``Localizing diagrams and deriving their Kan extensions'' at DER\_I-a, manuscript page 13, explains this distinction in general. What is new in these pages is the unusually economical categorical model. Total categories give $f_!$ without a separate cofibrant-replacement construction; the cost is paid on the inverse-image side by $\Phi_S$. The proof of the two triangles shows exactly why this economy works. In particular, the ordinary adjunction between postcomposition and pullback remains visible inside the homotopical adjunction, but its counit must be read through an invertible replacement map.

This geometric picture also explains why Grothendieck immediately turns to locally constant coefficients. Integration is useful both for arbitrary diagrams and for families whose transports are equivalences. The passage between those two kinds of coefficients is a localization of the coefficient theory itself. The next formal lemma separates that passage from the particular category-of-categories model, so that the same argument can subsequently be used wherever the relevant reflection exists.
\end{dcomment}

\begin{dcomment}{exposition}{p150:ce qui est la bijection annoncée.}{Reflecting coefficients while keeping the adjunction}
Grothendieck wants to transport the adjunction to locally constant objects without constructing its unit and counit a second time. The reflective localizations $\rho\dashv i$ and $\rho'\dashv i'$ are the right mechanism: maps into a local object already factor through the reflection. The lemma is entirely categorical, and its hypotheses should be separated from the work needed to produce these reflections in the homotopical situation.

Let $F:\mathcal A'\rightleftarrows\mathcal A:G$ be an adjunction and let $i:\mathcal A_0\hookrightarrow\mathcal A$ and $i':\mathcal A'_0\hookrightarrow\mathcal A'$ be fully faithful right adjoints to $\rho$ and $\rho'$. Assume that $G$ carries local objects to local objects. Its restriction is then $G_0:\mathcal A_0\to\mathcal A'_0$, characterized by $i'G_0\cong Gi$. A left adjoint can be written immediately as
\[
 F_0=\rho Fi'.
\]
Indeed, for $x\in\mathcal A'_0$ and $y\in\mathcal A_0$ there are natural bijections
\[
\begin{aligned}
 \Hom_{\mathcal A_0}(\rho Fi'x,y)
 &\cong\Hom_{\mathcal A}(Fi'x,iy)\\
 &\cong\Hom_{\mathcal A'}(i'x,Giy)\\
 &\cong\Hom_{\mathcal A'_0}(x,G_0y).
\end{aligned}
\]
This proves the adjunction directly. Each bijection comes from one of the given adjunctions or from full faithfulness, so its naturality, and consequently the triangle identities, are part of the construction.

One can also recover the descent square used in the text. Suppose $u:x\to x'$ is inverted by $\rho'$. For every $y\in\mathcal A_0$, the map
\[
 \Hom_{\mathcal A}(Fx',iy)\longrightarrow\Hom_{\mathcal A}(Fx,iy)
\]
is bijective: after adjunction it is precomposition by $u$ into the local object $Giy$. Applying $\rho\dashv i$ and then Yoneda shows that $\rho F(u)$ is invertible. Thus $F$ preserves the local equivalences and descends to the localization. In particular, applying this to the unit $x\to i'\rho'x$ gives
\[
 \rho Fx\xrightarrow{\sim}\rho Fi'\rho'x=F_0\rho'x.
\]
Under the present hypotheses, Grothendieck's condition (b) therefore supplies condition (a) as well. His proof assumes both and is correct; the observation explains why the two squares fit so naturally together.

For the application, $F=f_!$ and $G=f^*$. Once the local objects and their reflection have been established, restriction preserves local constancy, and the formula becomes
\[
 f_!^{\mathrm{lc}}\cong\rho_S f_!i_{S'}.
\]
It need not say that $f_!$ itself carries a locally constant object to a locally constant object. The reflection on the target is part of the formula. For example, with ordinary space-valued presheaves, extending the point at $0$ along $\{0\}\hookrightarrow[1]$ gives the diagram $\varnothing\to *$, which is not locally constant. Confusing the restricted inverse image with a restriction of both adjoints would lose precisely this example.

The later theory supplies the reflections systematically in its established settings. \citet[Proposition~3.9 and Theorem~3.12]{Cisinski2009LocallyConstant} proves that the inclusion of locally constant objects admits both adjoints. Holding the base $A$ fixed while varying an additional diagram shape gives a derivator of such coefficients. This does not assert that varying the base $A$ itself gives a derivator with the ordinary comma formulas. The formal lemma here explains the left-reflection part of the induced adjunction; it does not itself prove existence of a reflection for every tentative class $W$ considered in the manuscript.
\end{dcomment}

\begin{dcomment}{exposition}{p151:sont décrits p. 138 à 140.}{The product and the choice of representatives}
Grothendieck now wants a formula for integrating a restricted coefficient. Before such a formula can be meaningful, multiplication in $\HOT_W(S)$ must be constructed. The basic observation is that two smooth representatives may be multiplied over $S$, and that the resulting operation respects the localized homotopy types. The false corollary which follows arises from changing the quantifiers in this observation.

Let $\mathcal L_S$ be the category of smooth categories over $S$, with the colocal equivalences inherited from $\Cat/S$. Its inclusion induces the equivalence of localizations already used in the adjunction proof. Suppose $u:X\to X'$ lies in $W_S^{\mathrm d}$ and both $X$ and $X'$ are smooth. Then, for every $Y\to S$, the pullback
\[
 u\times_S Y:X\times_S Y\longrightarrow X'\times_S Y
\]
is a colocal equivalence over $Y$. Postcomposition with $Y\to S$ carries it to a colocal equivalence over $S$. This last passage is the compatibility established for $f_!$; it is not an assertion that arbitrary base change preserves arbitrary colocal equivalences.

Taking $Y$ smooth shows that the product bifunctor on $\mathcal L_S$ preserves weak equivalences separately in each variable. It therefore descends to the localization. A simultaneous change in both variables factors into two separate changes, so no further compatibility is needed. Transporting the descended functor across $\mathcal L_S[W^{-1}]\simeq\HOT_W(S)$ gives
\[
 Q_S(X)\times_S Q_S(Y)\cong Q_S(X\times_S Y)
 \qquad(X,Y\text{ smooth}).
\]
As usual, ``unique'' means unique with the specified compatible isomorphism on representatives, up to the corresponding unique natural isomorphism.

One can also see directly that this is a categorical product. The ordinary adjunction between the diagonal $\mathcal L_S\to\mathcal L_S\times\mathcal L_S$ and fibre product has both functors compatible with weak equivalences. Its unit, counit, and triangle identities consequently descend. Since localization of a product relative category is the product of the localizations, the descended fibre product is right adjoint to the diagonal. The object $Q_S(S)$ is terminal by the same argument. Associativity, symmetry, and the unit constraints are inherited coherently; they are not additional arbitrary choices of a monoidal structure.

The assertion needed in this proof concerns a weak equivalence whose \emph{two ends} are smooth. The attempted corollary instead applies it to $Y\to\Phi_S(Y)$ when $Y$ is arbitrary, on the ground that the fixed multiplier $X$ is smooth. That is a different assertion. Smoothness of the base-change functor does not make every colocal equivalence stable under its pullback. The interval example below proves that the difference persists after localization.

The boxed formula on manuscript page 151 must therefore be read as provisional. Its eventual unconditional form is $f_!f^*\cong\multhot_{S'/S}$. The product formula $f_!f^*(\xi)\cong Q_S(S')\times_S\xi$ is recovered when $S'\to S$ is smooth. It would be misleading to repair the proof of the false corollary while leaving that earlier boxed formula without its later qualification: the same issue affects both statements.
\end{dcomment}

\begin{dcomment}{examples}{p154:Le contrexemple le plus trivial}{The two ends of the interval}
Grothendieck's counterexample is particularly useful because every category involved has at most two objects. It reveals a distinction which is otherwise easy to conceal behind the word ``derived''. We work here with $W=W_\infty$, so that an empty homotopy type and a point are different, and put $S=[1]=(0\to1)$.

Let $X=\{0\}$ and $Y=\{1\}$, with their indicated maps to $S$. The map $X\to S$ is a Grothendieck fibration. Indeed, the only arrow of $S$ with target the image of its sole object is the identity of $0$, which lifts. Hence $X$ is smooth. The map $Y\to S$ is not a fibration: the arrow $0\to1$ has no lift. Their strict fibre product is empty:
\[
 X\times_S Y=\varnothing.
\]

The comma categories compute the replacement without any appeal to an abstract model structure. There is no arrow $1\to0$, while $0\to0$, $0\to1$, and $1\to1$ are unique. Consequently
\[
\begin{array}{c|cc}
 &s=0&s=1\\ \hline
 s\backslash X& e&\varnothing\\
 s\backslash Y& e&e
\end{array}
\]
and the transition for $Y$ is the identity of a point. More concretely, $\Phi_S(Y)=S/1\cong S$ over $S$, and $Y\to S$ is a colocal equivalence because both induced maps on comma categories are equivalences of contractible categories. Thus $Q_S(Y)$ is the terminal object of $\HOT_{W_\infty}(S)$.

The product of the two localized objects is therefore $Q_S(X)$, whereas the localization of their strict fibre product is initial:
\[
 Q_S(X)\times_S Q_S(Y)\cong Q_S(X),\qquad
 Q_S(X\times_S Y)=Q_S(\varnothing).
\]
They cannot be isomorphic. Evaluation at $0$ would turn such an isomorphism into an isomorphism between a point and the empty space. We only use the elementary fact that a functor preserves isomorphisms; no conservativity theorem for all evaluations is needed to obtain this contradiction.

The erroneous pullback step can be read in the same square. Pulling $Y\to\Phi_S(Y)=S$ back along the smooth map $X\to S$ gives
\[
 \varnothing=X\times_S Y\longrightarrow X\times_S S=X.
\]
It is not a colocal equivalence over $X$, since over the point this says that $\varnothing\to e$ is a weak homotopy equivalence. It is also not one over $S$, as the comma category at $0$ detects. This checks the exact implication used in the failed argument, rather than merely observing that its conclusion is implausible.

For orientation, a space-valued presheaf on $[1]$ is an arrow $F_1\to F_0$. Our three objects $Q_S(X)$, $Q_S(Y)$, and $Q_S(\varnothing)$ are respectively
\[
 (\varnothing\to *),\qquad(*\xrightarrow{\id}*),\qquad
 (\varnothing\to\varnothing).
\]
Their products are computed pointwise. This elementary description explains why the end $0$ is the smooth open end in Grothendieck's terminology and why choosing the other end for the coefficient destroys the naive strict-product formula.

The example is not meant to exclude every possible class $W$: if $\varnothing\to e$ has already been inverted, this particular distinction disappears. It disproves the unrestricted statement in its principal intended instance and in any nondegenerate instance in which the same empty-versus-point distinction survives. Grothendieck himself supplies the counterexample and the correction; the error is not an unresolved defect of the later theorem.
\end{dcomment}

\begin{dcomment}{exposition}{p155:ne doit pas être confondue}{Multiplication by a category over the base}
The corrected corollary retains an operation that is more informative than the product of two localized objects. Fix the category $X\to S$ itself and derive multiplication only in the other variable. Since a weak equivalence between two smooth representatives remains a weak equivalence after multiplication by any fixed $X$, the result is a well-defined endofunctor
\[
 M_X=\multhot_{X/S},\qquad
 M_X(Q_SY)=Q_S(X\times_S\Phi_S(Y)).
\]
The fixed category is part of the datum. The operation need not depend only on $Q_S(X)$.

The interval calculation makes that dependence completely explicit. Write a presheaf of spaces on $S=[1]$ as $F=(F_1\xrightarrow{u}F_0)$. The maps of bases $j_0:\{0\}\to S$ and $j_1:\{1\}\to S$ give
\[
\begin{aligned}
 M_{\{0\}}F&=(\varnothing\longrightarrow F_0),\\
 M_{\{1\}}F&=(F_1\xrightarrow{\id}F_1),\\
 M_SF&=(F_1\xrightarrow{u}F_0).
\end{aligned}
\]
To justify the first two formulas, restriction along $j_i$ evaluates at $i$. The left Kan extension of a space $T$ from $0$ is $\varnothing\to T$: there is no arrow $1\to0$ in $S$. From $1$ it is $T\xrightarrow{\id}T$: there is one arrow from each object of $S$ to $1$. In these cases the indexing categories are empty or terminal, so the formulas are already homotopy-correct. Corollary~2 identifies these composite functors with $M_{\{i\}}$.

Although $Q_S(\{1\})\cong Q_S(S)$, the last two endofunctors are not isomorphic. Taking $F=Q_S(\{0\})$ gives the empty diagram for $M_{\{1\}}F$, and the nonempty diagram $F$ for $M_SF$. Thus there cannot be a bifunctor on localized objects whose value at $(Q_SX,Q_SY)$ is always $M_X(Q_SY)$. The first-variable ambiguity is an actual loss of information, not a defect in the notation.

The adjunction also gives $M_X$ a counit and a comultiplication, because
\[
 M_X\cong(p_X)_!p_X^*,\qquad p_X:X\to S.
\]
Its counit is the map described by the projection in the preceding proof; its comultiplication is $(p_X)_!\alpha p_X^*$. The two comonad identities and coassociativity follow from the two triangle identities and naturality of $\alpha$. For $j_1$, the counit at $F$ has components $\id_{F_1}$ and $u:F_1\to F_0$. It is an isomorphism precisely when $F$ is locally constant. This gives an especially small illustration of the link between the endofunctor $f_!f^*$ and the behavior of transport maps.

When $X$ is smooth, the fixed representative is itself admissible in the product construction and $M_X$ is multiplication by $Q_S(X)$. Applying this with $X=S'$ gives the qualified product formula in Corollary~2. One can push the same calculation slightly further without introducing a new hypothesis: if $f:S'\to S$ is smooth, then
\[
 f_!(A\times_{S'}f^*B)\cong f_!A\times_SB.
\]
Choose smooth representatives $Z\to S'$ for $A$ and $Y\to S$ for $B$. The left side is represented over $S$ by $Z\times_{S'}(S'\times_SY)=Z\times_SY$. Since smooth maps compose, $Z\to S$ is smooth, so this also represents the right side. The isomorphism is natural on these representatives and descends to the localization. The formula for $f_!f^*$ is its case $A=1_{S'}$. This extension is a direct consequence of the same argument, not a separate assertion of originality.
\end{dcomment}

\begin{dcomment}{exposition}{p156:f_!f^*(\xi)\simeq\hot}{The comma square contains a natural transformation}
Grothendieck is now turning an existence theorem for left adjoints into a calculation of their values. The comma category has exactly the required information: it records not only objects over the chosen point but also arrows from that point to their images. These arrows are why an ordinary fibre cannot replace it.

Let $f:S'\to S$ and $s\in S$. Put $K=(s\downarrow f)$, with projection $q:K\to S'$ and structural functor $r:K\to e$. Its objects are pairs $(a,\lambda:s\to f(a))$, and its morphisms are maps $v:a\to a'$ satisfying $f(v)\lambda=\lambda'$. There is a specified natural transformation
\[
 sr\Longrightarrow fq,
\]
whose component is $\lambda$. Thus the first square on manuscript page 157 is a \emph{comma square}. The two composites are usually unequal, and the displayed arrows $\lambda$ need not be invertible. The word ``commutatif'' in the text should be read with this additional two-dimensional structure supplied, not as an assertion of strict commutativity.

In the presheaf convention, that natural transformation induces $q^*f^*\to r^*s^*$. Taking its left-adjoint mate gives the canonical comparison
\[
 r_!q^*\longrightarrow s^*f_!.
\]
The edition writes the arrow in the opposite direction on page 157; once the comparison is proved invertible its inverse has that direction. The comparison in Theorem~4, page 162, has the conventional direction. Keeping track of the mate here avoids defining an isomorphism merely because its source and target have happened to receive the same description.

For a representative $X'\to S'$, the right side has the concrete value
\[
 s^*f_!Q_{S'}(X')=Q_e(s\backslash X').
\]
The other side gives the same value without changing the representative. Indeed, $q$ is a discrete opfibration, and its categorical pullback of $X'$ is canonically $s\backslash X'$ over $K$. There is a direct way to check that this pullback computes the derived inverse image. For $k=(a,\lambda)\in K$, an object of $k\backslash q^*X'$ consists of $x\in X'$ and a morphism $v:a\to p(x)$; the remaining arrow in $K$ is forced to be $f(v)\lambda$. This gives a canonical isomorphism of categories
\[
 k\backslash(q^*X')\cong a\backslash X'.
\]
It is compatible with maps of representatives. Consequently the raw pullback functor $q^*$ carries every colocal equivalence to a colocal equivalence. It may therefore be used on all of $\Cat/S'$, and
\[
 r_!q^*Q_{S'}(X')=Q_e(s\backslash X').
\]
Under these constructions the mate is the identity map of this comma-category model. This proves the required pointwise formula, including its canonical comparison, rather than just an objectwise coincidence.

For $W_\infty$ the modern presheaf expression is
\[
 (f_!F)(s)\simeq
 \operatorname*{hocolim}_{(s\downarrow f)^{\mathrm{op}}}F\circ q^{\mathrm{op}}.
\]
The opposite on the indexing category is required because $F$ is contravariant. In the categorical model, $s\backslash X'$ is the total category representing this homotopy colimit. The familiar pointwise Kan formula has thus emerged from postcomposition of total categories and the single replacement $\Phi$. The corresponding general derivator statements are the exactness of comma squares and its equivalence with the pointwise axiom; see \citet[Proposition~1.26]{Groth2013}, after reversing the indexing convention. The detailed mate calculus was introduced in the earlier companion at DER\_I-a, manuscript page 31; here the comma category supplies a particular exact square to which that calculus applies.
\end{dcomment}

\begin{dcomment}{examples}{p158:p^*f_!\simeq f'_!q^*.}{Induction along a homomorphism of groups}
The comma formula also explains a familiar operation with monodromy. It gives a useful example in which the strict fibre can be connected while the derived fibre has many components. Let $\varphi:H\to G$ be a homomorphism of discrete groups and let $f:BH\to BG$ be the functor between their one-object groupoids. We again take $W=W_\infty$.

At the unique object $s$ of $BG$, the comma category $K=(s\downarrow f)$ has one object for each $g\in G$. A morphism from $g$ to $g'$ is an $h\in H$ with
\[
 \varphi(h)g=g'.
\]
Composition is multiplication in $H$. Thus $K$ is the action groupoid for the left action of $H$ on $G$ through $\varphi$. Its connected components are indexed by $\varphi(H)\backslash G$, and the automorphism group of every object is $\ker\varphi$. Choosing one object in each orbit gives an equivalence of groupoids
\[
 K\simeq\coprod_{\varphi(H)\backslash G}B(\ker\varphi).
\]
The choice is unnecessary for the homotopy type; it only writes that type as an explicit disjoint union.

Apply the pointwise formula to the terminal coefficient on $BH$. It is represented by the identity category $BH\to BH$, so its left direct image is represented by $BH\to BG$. Consequently
\[
 (f_!1)(s)\simeq
 \coprod_{\varphi(H)\backslash G}B(\ker\varphi).
\]
By contrast, the strict fibre of $BH\to BG$ has one object and automorphism group $\ker\varphi$. It sees one component, regardless of the index of $\varphi(H)$ in $G$. For a proper subgroup inclusion $H\hookrightarrow G$, the strict fibre is a point and the derived fibre is the discrete coset set $H\backslash G$. For a surjection, there is one component and both descriptions have homotopy type $B(\ker\varphi)$. In that case the functor is a Grothendieck fibration: every arrow of $BG$ has a lift, and every arrow between groupoids is cartesian. These calculations exhibit exactly the distinction expressed by smoothness in the general formula.

Right multiplication on $G$ commutes with the left $H$-action. It therefore acts on $K$, describing the contravariant transport of $f_!1$ around $BG$. More generally, contravariant diagrams on $BH$ are right $H$-objects, with the appropriate coherent homotopical interpretation. For an inclusion $H\subseteq G$ and a right $H$-space $T$, the left direct image is ordinary induction
\[
 T\times_HG,
 \qquad (th,g)\sim(t,hg),
\]
with its right $G$-action. Choosing representatives for $H\backslash G$ expresses the underlying space as a disjoint union of copies of $T$, so induction preserves weak equivalences and no further derivation changes this model. Its unit sends $t$ to $[t,1]$, and the counit at a right $G$-space $Z$ sends $[z,g]$ to $zg$; the balancing relation makes the latter map well defined. The identities $[t,1]g=[t,g]$ and $z1=z$ verify the triangles in this example.

For a noninjective homomorphism the balanced product must instead be derived. The terminal coefficient already shows why: the ordinary orbit set would forget the factors $B(\ker\varphi)$. Grothendieck's comma-category construction retains them automatically. Thus the same pointwise formula accounts both for induction of group actions and for the homotopy fibre information carried by a general map of classifying spaces.
\end{dcomment}

\begin{dcomment}{status}{p159:La construction des $f_*$}{What Theorem 3 establishes at this point}
Grothendieck is assembling the axioms of a derivator, and the annotation to Theorem~3 is part of its mathematical content. The left adjoints and their comma-category formula have been proved. The assertion concerning detection of isomorphisms by fibres, called Der~2 here, is expressly withdrawn in footnote~76. The theorem must be read with that qualification, rather than as an unconditional construction of a full derivator followed by an inessential caution.

The proof of the constant-coefficient assertion is complete and especially short. If $p:S\to e$, $p':S'\to e$, and $f:S'\to S$ belongs to $W$, then for $\xi\in\Hot_W$ the comparison
\[
 p'_!p'^*\xi\longrightarrow p_!p^*\xi
\]
is multiplication of $Q_e(f)$ by $\xi$. In this formula the products take place over $e$:
\[
 Q_e(S')\times\xi\longrightarrow Q_e(S)\times\xi.
\]
The subscripts $S$ printed in the product on manuscript page 158 are a notational slip. Since $Q_e(f)$ is invertible, the displayed map is invertible. This proves the stated constant-coefficient invariance, without any use of Der~2.

The missing axiom concerns a different issue. Even if weak equivalences of representatives were defined by inspecting comma categories, a general morphism in the localized category is not necessarily represented by one original functor. Detecting its invertibility pointwise therefore requires control of localization itself. Strong saturation of $W$ helps identify those original arrows which become invertible over the point; by itself that identification does not constitute a proof for every morphism in every diagram category. The manuscript returns to precisely this distinction on pages 165--166.

The modern theory resolves this issue for fundamental localizers. \citet[Propositions~3.1.30 and 4.2.7]{Cisinski2006} proves pointwise detection for presheaf localizers and the diagramwise comparison with categories; taking the simplex test category transfers the result to the present categorical presentation. The commentary at manuscript page 165 explains that repair. This does not change the status of the argument here: the required theorem was still missing at this point in the manuscript. In the standard space case, a familiar alternative is the strong derivator of a combinatorial model category, established in \citet[Proposition~1.30]{Groth2013}; it provides both homotopy Kan extensions as well. The manuscript's replacement calculus and these later localization theorems perform different parts of the construction.

The historical numbering also needs care. The manuscript's Der~4$'$ is existence of the left adjoints, and Der~5$'$ is their pointwise calculation. In the modern convention of the cited paper, these belong respectively to (Der3) and (Der4). Translating a theorem number without translating its statement would make the audit look stronger than the argument actually is. The prospective internal-Hom proof mentioned at the end of the page is not needed for the stronger locally constant assertion: the following unit-and-counit calculation gives a direct route.
\end{dcomment}

\begin{dcomment}{later}{p161:qui est parfait, est dans}{Why local systems follow the homotopy type of the base}
Grothendieck's next goal is stronger than invariance of constant-coefficient homology. He wants a weak equivalence of bases to identify the whole categories of locally constant coefficients. The proof here is an application of the adjunction already constructed, but it uses an additional stability hypothesis on $W$. The hypothesis has real content and should remain visible in the modern formulation.

Assume the preceding theory of locally constant replacements and the axiom L6: pullback of a $W$-equivalence along a perfect functor is again a $W$-equivalence. Take $f:S'\to S$ in $W$. For a locally constant coefficient represented by a perfect map $X\to S$, the counit of $f_!^{\mathrm{lc}}\dashv f^*_{\mathrm{lc}}$ is represented by
\[
 X\times_SS'\longrightarrow X.
\]
This is the image of the previously computed counit under reflection, so its identification with the projection is justified by the localized-adjunction lemma. L6 makes it a $W$-equivalence and hence an isomorphism in the coarse localization.

For the unit, take $X'\to S'$ and replace its composite over $S$ by the locally constant replacement $C_S(X')$. There is a canonical map
\[
 a:X'\longrightarrow C_S(X')\times_SS'
\]
whose first component is $i_{X'}$ and whose second is the structure map. The projection $\pi:C_S(X')\times_SS'\to C_S(X')$ is a $W$-equivalence by L6 because $C_S(X')\to S$ is perfect, and $\pi a=i_{X'}$ is a $W$-equivalence by the replacement construction. Two-out-of-three gives $a\in W$. Thus both unit and counit are isomorphisms and
\[
 f^*_{\mathrm{lc}}:\HOT_W^{\mathrm{lc}}(S)
 \xrightarrow{\sim}\HOT_W^{\mathrm{lc}}(S')
\]
has quasi-inverse $f_!^{\mathrm{lc}}$. The symbol $\beta_{0,X'}$ used for this unit on page 161 is a slip of notation: its direction is that of $\alpha_0$.

There is a precise converse behind the footnote on necessity. Suppose restriction along every $f\in W$ is fully faithful on locally constant objects and that the displayed projection formula for the counit is available on perfect representatives. A right adjoint is fully faithful exactly when its counit is invertible. The projection $X\times_SS'\to X$ therefore becomes invertible in the coarse localization, and hence after forgetting the base in $\Hot_W$. If $W$ is strongly saturated, the projection itself belongs to $W$. This recovers L6. The stronger L6~bis requires the same argument for homotopy-fibred representatives and their raw pullback description; it is not obtained merely by deleting the word ``perfect'' from this proof.

In the standard theory of spaces the homotopy invariance is a theorem with much greater generality in the coefficients. \citet[Theorem~1.3]{Cisinski2009LocallyConstant} proves that a functor inducing a weak equivalence of nerves induces an equivalence on the categories of locally constant diagrams with values in a model category. Applied after taking opposites, this gives the variance used here. It is a later precise form of the assertion under discussion, not a claim that the manuscript has already proved every required replacement theorem for an arbitrary $W$.

The interval gives a compact interpretation. A locally constant presheaf on $[1]$ is an equivalence $F_1\to F_0$ of spaces. Evaluation at either end is an equivalence from its homotopy category to that of spaces: a quasi-inverse sends $T$ to the constant arrow $T\xrightarrow{\id}T$, and the given equivalence supplies the comparison back to $F$. The map from either endpoint to $[1]$ is a weak equivalence of nerves. The locally constant left adjoint includes the reflection discussed above, which is why this conclusion is compatible with the different unrestricted left Kan extensions from the two endpoints.
\end{dcomment}

% Commentary for manuscript pp. 162--184. See findings.md for source checks.
\begin{dcomment}{exposition}{p162}{Why the two base-change hypotheses are different}
Grothendieck wants the left direct image to behave as an integration along the fibres of a category. The preceding adjunction gives such an operation, but does not yet say when integration commutes with a change of the base. The square on this page supplies the comparison; smoothness of its vertical side and properness of its horizontal side are two different reasons why that comparison is invertible. The terminology is inherited from the cohomological formalism discussed in the commentary to Chapter~I, manuscript pp.~31 and 52--58. Here we can see the categorical mechanism directly, before using a general theorem about derivators.

Write $Q_T(Z)=\hot_{T,W}(Z)$ and retain the square with $X'=X\times_S S'$. The natural transformation in question has direction
\[
 f'_!q^*\longrightarrow p^*f_!.
\]
It is the composite obtained by applying $f'_!q^*$ to the unit of $f_!\dashv f^*$, using $q^*f^*\simeq f'^*p^*$, and then applying the counit of $f'_!\dashv f'^*$. This description matters: identifying the values of two functors does not, by itself, identify the particular comparison map between them.

Suppose first that $f$ is smooth. Represent $\xi\in\HOT_W(X)$ by $Z\to X$ smooth over $X$. Composition makes $Z\to S$ smooth as well. Both restrictions can therefore be computed by ordinary fibre products, whereas both left direct images are computed by composition of structure maps. Consequently
\[
 f'_!q^*Q_X(Z)=Q_{S'}(Z\times_S S')=p^*f_!Q_X(Z).
\]
With these representatives the mate above is the identity of the displayed object. One can check this without any abstract exactness criterion: the unit inserts $Z$ into its comma-category replacement, and the counit projects from the corresponding fibre product; their composite over $S'$ is the comparison induced by the identity of $Z\times_S S'$. The unit and counit calculated on pp.~138--146 therefore supply the required compatibility. This proves the theorem for smooth $f$.

The proper case changes a different part of the calculation. Here $f_!Q_X(Z)=Q_S(Z)$ need not have a smooth representative of this particular form. What is needed is that ordinary pullback by $p$ preserve \emph{all} colocally weak equivalences in $\Cat/S$. If this is known, the canonical map
\[
 Z\times_S S'\longrightarrow\Phi_S(Z)\times_S S'
\]
is a colocally weak equivalence, so ordinary pullback already computes the derived restriction on every $Z$. We can still choose $Z$ smooth over $X$ to calculate $q^*$, obtaining the same identity as before. Thus the two hypotheses solve distinct problems: the first makes the chosen representative suitable over $S$; the second makes every representative suitable for this particular change of base.

The modern characterization of smooth functors by these universal base-change identities is given in \citet[Theorem~1.2]{Maltsiniotis2005Asphericite}. His treatment also identifies the concrete comparison with the adjunction mate, in Remark~5.5. These results retain the direction of the comma category used here: $s\bs X$ consists of arrows $s\to f(x)$. Replacing it by $X/s$ changes the variance and exchanges the relevant smooth and proper conditions.
\end{dcomment}

\begin{dcomment}{exposition}{p165}{What the constant coefficients actually detect}
Grothendieck pauses the base-change argument to identify the weak equivalences seen by his new coefficients. This is useful because properness can then be compared with the corresponding notion for a derivator. Proposition~5 is also a precise safeguard against confusing an invertible morphism in a localization with membership in the class originally chosen for inversion.

Let $H=\Hot_W$ and let $a_S=p_S^*a$ denote a constant coefficient. The adjunction and the product calculation give a natural bijection
\[
 \Hom_{\HOT_W(S)}(a_S,b_S)
 \simeq\Hom_H\bigl(Q_e(S)\times a,b\bigr).
\]
For $f:S'\to S$, the restriction map on the left corresponds on the right to precomposition with $Q_e(f)\times\id_a$. Thus restriction induces bijections for every $a,b$ precisely when $Q_e(f)\times\id_a$ is invertible for every $a$. One implication uses Yoneda in $H$; for the other, an isomorphism remains an isomorphism under a functor. Taking $a$ to be the unit $Q_e(e)$ shows that it is enough to test $Q_e(f)$ itself. This proves Proposition~5 without a criterion detecting isomorphisms of arbitrary relative objects by their fibres.

The displayed Hom groups immediately before the manuscript marker~165 have incorrect subscripts in the edition. After applying $p_{S!}$ and $p_{S'!}$, their arguments lie in $H$, so both Hom groups belong to $H$, not to $\HOT_W(S)$ and $\HOT_W(S')$. The types of the functors determine the correction uniquely. We preserve the printed formula in the transcription and record the slip in the corrections table; the argument just given uses the correctly typed formula.

Write $\overline W$ for the class of arrows inverted by $Q_e:\Cat\to H$. What has been proved is
\[
 \{\HOT_W\text{-equivalences}\}=\overline W.
\]
The extra hypothesis of strong saturation identifies $\overline W$ with $W$. This is a statement about the original arrows of $\Cat$. It should be distinguished from the assertion that the evaluation functors $\HOT_W(S)\to H$ jointly detect invertibility of \emph{every} morphism of a localized relative category. Such a morphism may be represented by a zigzag rather than by a single functor. Strong saturation of $W$ does not provide the missing passage from that zigzag to a pointwise invertibility argument. This is exactly the unfinished step Grothendieck identifies on p.~166; he does not overlook it. Under the modern fundamental-localizer hypotheses, strong saturation itself follows from \citet[Proposition~4.2.4]{Cisinski2006}; that later theorem does not turn the argument printed here into a proof of the general fibrewise criterion.

The gap was subsequently resolved for every fundamental localizer, without an accessibility assumption. The split comma-category presentation identifies $\HOT_W(S)$ with the localization of $\Cat^{S^{\mathrm{op}}}$ at pointwise $W$-equivalences: its two comparisons are $X\to\Phi_S(X)$ and the transport retraction from the comma replacement of a split fibration. Choose the simplex category as a $W$-test category. \citet[Proposition~4.2.7]{Cisinski2006} then compares these diagram localizations with the corresponding localizations of simplicial-set diagrams; Proposition~3.1.30 establishes joint conservativity of their point evaluations for an arbitrary presheaf localizer. The comparison commutes with evaluation, so the criterion transfers to $\HOT_W(S)$. This is a general later repair of Der~2 under the revised W(1,2,3), rather than a restriction to accessible examples. In the manuscript, the direct proof on pp.~168--169 avoids waiting for it.
\end{dcomment}

\begin{dcomment}{development}{p167}{Testing base change on a single coefficient}
Grothendieck tests a relative equivalence by homology with every constant coefficient. The terminal coefficient already gives a useful sharpening of his argument. It isolates the precise amount of base change needed to recover preservation of the relative localizer, and explains why the unfinished proof of the fibrewise isomorphism criterion is irrelevant to this particular converse.

\textbf{Terminal-coefficient criterion.}\ours\
Assume the constructions of $\HOT_W$, its restrictions, and its left direct images established before p.~162, with $W$ strongly saturated. Fix $p:S'\to S$. For every $f:X\to S$, put $X'=X\times_S S'$ and write $q:X'\to X$, $f':X'\to S'$. The following conditions are equivalent:
\begin{enumerate}
\item Ordinary pullback $p^{-1}:\Cat/S\to\Cat/S'$ takes $W_S^{\mathrm d}$ into $W_{S'}^{\mathrm d}$.
\item For every $f$ and every $\xi\in\HOT_W(X)$, the base-change mate $f'_!q^*\xi\to p^*f_!\xi$ is invertible.
\item For every $f$, that mate is invertible on the terminal coefficient $\mathbf1_X$ alone.
\end{enumerate}
Here $\mathbf1_X$ is represented by $X\xrightarrow{\id}X$; it is the constant coefficient with value the terminal homotopy type. Strong saturation is used only to recover membership in $W$ from an isomorphism in $\Hot_W$.

\textit{Proof.} We first record a detection statement for actual functors. If $u:Z\to T$ is over a category $B$, then
\[
 u\in W_B^{\mathrm d}\quad\Longleftrightarrow\quad
 Q_B(u)\text{ is invertible in }\HOT_W(B).
\]
The forward implication is the definition of localization. For the converse, evaluate the invertible morphism at each $b\in B$. The established pointwise formula identifies its image with $Q_e(b\bs u)$. Strong saturation gives $b\bs u\in W$ for all $b$, which is exactly $u\in W_B^{\mathrm d}$. We have used preservation of isomorphisms by evaluation, not reflection of isomorphisms by the family of evaluations.

Assume condition~1. The maps $Z\to\Phi_S(Z)$ are in $W_S^{\mathrm d}$, so their pullbacks are inverted by $Q_{S'}$. They identify the functor induced by ordinary pullback with $p^*$. Given $\xi$, choose a smooth representative $Z\to X$. Then $q^*\xi$ is represented by $Z\times_S S'$, and condition~1 identifies $p^*f_!\xi$ with the same object over $S'$. The unit--counit calculation of the preceding exposition identifies the mate with this canonical comparison. Thus condition~2 holds, and condition~2 immediately implies condition~3.

For the converse, observe that
\[
 f_!\mathbf1_X=Q_S(X),\qquad
 f'_!q^*\mathbf1_X=Q_{S'}(X\times_S S').
\]
The mates in condition~3 therefore give a natural isomorphism, for $X\in\Cat/S$,
\[
 \theta_X:Q_{S'}(p^{-1}X)\xrightarrow{\sim}p^*Q_S(X).
\]
For a functor $u:X\to Y$ over $S$, naturality gives
\[
 \theta_Y\,Q_{S'}(p^{-1}u)=p^*Q_S(u)\,\theta_X.
\]
If $u\in W_S^{\mathrm d}$, the right-hand side is a composite of isomorphisms. Hence $Q_{S'}(p^{-1}u)$ is invertible. The detection statement already proved, now over $S'$, gives $p^{-1}u\in W_{S'}^{\mathrm d}$. This is condition~1.\hfill$\square$

This is a development of the argument actually present on pp.~166--167. We checked the universal base-change formulations in \citet[Theorem~1.2 and \S5]{Maltsiniotis2005Asphericite}; we did not find this exact three-condition formulation isolated there. The assertion is a formal consequence of the manuscript's construction, and the label records our presentation and proof, without a claim of priority. Its practical value is that one may test a prospective base-change theorem on the objects $Q_S(X)$ before attempting a theorem for arbitrary coefficients.
\end{dcomment}

\begin{dcomment}{exposition}{p169}{The direct argument and the missing manuscript page}
Grothendieck returns to the categories themselves because the obstruction was caused by a premature appeal to the general derivator formalism. The direct proof is short once the appropriate comma categories have been chosen. Its force is that it proves the proper case of Theorem~4 before the general fibrewise criterion has been established.

The elementary observation is valid for any class $W$ with the two-out-of-three property. Call $g:Z'\to Z$ a universal $W$-equivalence if every projection $T\times_Z Z'\to T$ belongs to $W$. For a functor $u:T\to U$ over $Z$, the pullback square gives an equality
\[
 u\,g_T=g_U\,g^*u.
\]
Both $g_T$ and $g_U$ lie in $W$. Consequently $u$ belongs to $W$ if and only if $g^*u$ does. Thus pullback along a universal $W$-equivalence both preserves and reflects total $W$-equivalences. No saturation beyond two-out-of-three is involved.

For a proper functor $p:S'\to S$, fix $s'\in S'$ and put $s=p(s')$. The relevant functor is
\[
 r_{s'}:s'\bs S'\longrightarrow s\bs S.
\]
The categorical characterization of properness says that $r_{s'}$ is a universal $W$-equivalence; in the original setting one first uses the usual homotopy class and then its inclusion in $W$. The fibre of $r_{s'}$ over an arrow $s\to t$ is the category of its lifts starting at $s'$. Properness provides the required asphericity of these lift categories, together with the stability needed after further base change. This is the same lifting-category mechanism explained in the commentary to Chapter~I, p.~52, now used for a specific computation.

Suppose $u:X\to Y$ is in $W_S^{\mathrm d}$. After base change to $s\bs S$, its total functor belongs to $W$ by definition. Pulling this functor back along $r_{s'}$ gives precisely
\[
 s'\bs(X\times_S S')\longrightarrow s'\bs(Y\times_S S').
\]
The elementary observation proves that this functor is in $W$. Since $s'$ was arbitrary, $p^{-1}u\in W_{S'}^{\mathrm d}$. The calculation on p.~163 now completes proper base change. The earlier detour through Der~2 is unnecessary for this theorem, as Grothendieck's footnote explicitly acknowledges; it remains a genuine unfinished task for the general theory.

One should retain the distinction between two conditions on a functor. A universal $W$-equivalence must itself be in $W$, by taking $T=Z$. A proper functor need not be an equivalence. It is the induced functor between the two comma categories, not the original $p$, to which the universal-equivalence lemma is applied. This distinction prevents a common but incorrect shortening of the proof.

The edition then prints ``page 170 manque''. There is no mathematical text to reconstruct at that point in the supplied source. The proof just completed reaches its stated conclusion on p.~169; the next surviving page begins a reorganization of the hypotheses. The absence of p.~170 must therefore be recorded as a source lacuna, not filled by a conjectural transition or counted as an additional gap in Theorem~4.
\end{dcomment}

\begin{dcomment}{status}{p172}{The minimum is the ordinary homotopy localizer}
Grothendieck now asks how much of the construction belongs to the ordinary homotopy theory of spaces, and how much follows from a few categorical principles. The distinction is productive even though his hope for a genuinely smaller fundamental localizer was later disproved. It separates the locality principle used to construct the theory from the stronger principle that detects an equivalence on each fibre.

Here W(1) comprises identities, two-out-of-three, and the condition that $fg,gf\in W$ imply $f,g\in W$. W(2) is the relative theorem~A condition: for $u:X\to Y$ over $S$, membership of every $X/s\to Y/s$ in $W$ implies $u\in W$. Its weaker form W(2~bis) only tests $X/y\to e$ for a functor $X\to Y$. W(3) declares a category with a terminal object to be aspherical. In particular, W(1) implies the retract clause used in the modern definition of a weak basic localizer: if $ri=\id$ and $ir\in W$, then $i,r\in W$. Thus the later minimum theorem applies to the weaker triple W(1,2~bis,3), not merely to the stronger relative version.

\citet[Theorem~2.2.11, p.~136]{Cisinski2004Minimal} proves that the weak homotopy equivalences of nerves form the smallest fundamental localizer. The stronger weak-localizer formulation is stated explicitly in \citet[Theorem~1.5]{Cisinski2009LocallyConstant}, with proof in \citet[Theorem~6.1.18]{Cisinski2006}. Consequently the class called $W_\omega$ here is $W_\infty$: the reverse inclusion follows because $W_\infty$ itself satisfies the stated axioms. The proposed homotopy theories with $W\subsetneq W_\infty$ therefore do not occur under these hypotheses. This is a settled question, rather than a still-open direction in the foundations of derivators.

That answer does not make the remaining distinctions disappear. A class $W$ may properly \emph{contain} $W_\infty$. It then identifies more maps than ordinary weak homotopy theory does. Rational homology equivalences give one such example. Nothing in the minimum theorem says that a total $W$-equivalence between relative objects must be a $W$-equivalence on every fibre. That is the additional assertion W(4b). The example below at p.~184 shows its failure even for genuine bifibrations with connected fibres.

The notation also needs historical care. $\Fib_W S$ in the revised argument denotes the weak $W$-fibrations of Chapter~XII, not the fibrations of a chosen Quillen model structure. Nor is it the earlier category $\Fib S$ of Grothendieck fibrations. One useful special case is unambiguous: a Grothendieck fibration belongs to $\Fib_W S$ exactly when every transport functor between its fibres lies in $W$. This is Chapter~XII, Theorem~3 and Corollary~2 on manuscript pp.~51--52 \citep{GrothendieckDerivateursXII}. A genuine bifibration with equivalence transports consequently supplies examples for every $W$ under consideration.

The names $W_\omega$ and $W'_\omega$ should remain in the historical text, because they record which hypotheses each construction was intended to use. Modern identifications are statements about those definitions, not reasons to silently replace their symbols throughout the manuscript. The continuing mathematical issue is how localization interacts with transport and fibre detection after $W$ has been enlarged.
\end{dcomment}

\begin{dcomment}{exposition}{p174}{The five tests and their domains}
Grothendieck organizes the relative theory by asking what it means for a functor over $S$ to be an equivalence. The same underlying functor can pass one test and fail another. A good class of representatives is useful precisely because several tests then become equivalent; the large diagram on p.~175 records choices of such representatives.

For $u:X\to Y$ over $S$, the tests are
\[
\begin{array}{ll}
 u\in W_S & \text{if the total functor }u\text{ belongs to }W,\\
 u\in W_S^{\mathrm d} & \text{if }s\bs X\to s\bs Y\text{ belongs to }W\text{ for every }s,\\
 u\in W_S^{\mathrm g} & \text{if }X/s\to Y/s\text{ belongs to }W\text{ for every }s,\\
 u\in W_S^{\mathrm f} & \text{if }X_s\to Y_s\text{ belongs to }W\text{ for every }s,\\
 u\in W_S^{\mathrm{univ}} & \text{if every pullback of }u\text{ along }T\to S\text{ belongs to }W.
\end{array}
\]
The last test concerns changes of the \emph{base} $S$ of the relative objects. These definitions are meaningful on $\Cat/S$, but the strict-fibre test does not by itself provide a satisfactory presentation of the desired relative homotopy theory there. This is the force of the warning that $W_S^{\mathrm f}$ should not be used on the whole slice.

To see the issue completely in the smallest useful example, take $S=[1]=(0\to1)$ and $W=W_\infty$. Let $u:\{1\}\to S$ be the inclusion, viewed as a morphism in $\Cat/S$ from the object supported at~1 to the terminal object $S\to S$. For $s=0$, the category $s\bs\{1\}$ is a point, and $s\bs S$ is $[1]$; for $s=1$, both are points. Thus $u\in W_S^{\mathrm d}$. Its strict fibre at~0 is $\varnothing\to e$, so $u\notin W_S^{\mathrm f}$. Its incoming comma map at~0 is also $\varnothing\to e$, so $u\notin W_S^{\mathrm g}$. Both total categories are contractible, and therefore $u\in W_S$. Dually, $\{0\}\to S$ is in $W_S^{\mathrm g}\cap W_S$ but not in $W_S^{\mathrm d}$. These computations separate the three localizers before any replacement is made.

For smooth objects, the natural comparisons $X_s\to s\bs X$ identify the strict-fibre and outgoing-comma tests. The locality theorem then promotes the fibre test to a total equivalence after every change of base. Conversely, a universal equivalence can be tested after $e\xrightarrow{s}S$, so it passes the fibre test. This explains
\[
 W_S^{\mathrm f}=W_S^{\mathrm d}=W_S^{\mathrm{univ}}
 \quad\text{on }\ul{Liss}_W S.
\]
It also explains the implication to $W_S^{\mathrm g}$: a universal equivalence remains an equivalence after $S/s\to S$. The proper case is the dual argument. On weak $W$-fibrations, both comma comparisons are available, and all four relative tests agree. They still imply the total test, but the converse is exactly W(4b).

The phrase ``this implication is strict if we do not impose W(4b)'' should be read as a statement about what the weaker assumptions guarantee. It does not mean that every $W$ for which the axiom has merely been left unstated yields strict inclusion. For $W_\infty$ the converse holds. What the next construction needs is a hypothesis that supplies that converse uniformly, and the rational example shows why it cannot be omitted from a theorem formulated for arbitrary $W$.
\end{dcomment}

\begin{dcomment}{examples}{p176}{A complete local system over a loop}
Grothendieck's upper floor consists of relative objects whose transport is invertible in the homotopy theory. An elementary calculation shows why ``locally constant'' should retain the entire system of transport, rather than just a common isomorphism class of fibres.

Let $S=B\mathbb Z$ be the category with one object and automorphism group $\mathbb Z$. Choose a set $E$ and a permutation $T:E\to E$. The Grothendieck construction of this action is a category $X$ over $S$: its objects are the elements of $E$, and an arrow $x\to y$ over $n\in\mathbb Z$ exists exactly when $y=T^n x$, in which case it is unique. Composition adds the integers. Every arrow is invertible, and a specified base arrow has a unique lift with prescribed source or prescribed target. Thus $X\to S$ is both a discrete fibration and a discrete opfibration. Its fibre is the discrete category $E$, and transport around the generator is $T$ or $T^{-1}$ according to the chosen variance. Since both are isomorphisms, the object lies in the locally constant part for every localizer used here.

Take $E=\{a,b\}$. When $T$ is the identity, $X$ is the disjoint union of two copies of $B\mathbb Z$. When $T$ exchanges $a$ and $b$, the category $X$ is connected: there is an arrow over~1 from $a$ to $b$. Its automorphism group at $a$ is $2\mathbb Z$, so the inclusion of the full one-object subcategory at $a$ is an equivalence of categories. Its nerve therefore has the homotopy type of a circle, whereas the first example has two circle components. The fibres of the two examples are the same two-point set; the relative objects and even their total homotopy types are different.

This example also makes the two localizations on the upper floor concrete. The fine localization identifies maps which are equivalences on the displayed fibres. The coarse localization identifies maps which are equivalences after taking the total category. For $W_\infty$, a map between these locally constant systems over the \emph{same} base is a total equivalence exactly when it is a fibre equivalence. For larger $W$ that conclusion can fail, because total invariants can lose the information carried by monodromy. The homology example at p.~184 exhibits that loss with connected fibres.

The later theory of locally constant functors gives an intrinsic setting for this picture. \citet[Theorem~3.12]{Cisinski2009LocallyConstant} constructs their derivator and proves that its inclusion into the shifted derivator has both adjoints. His formula~(4.1.2) identifies the ordinary-homotopy coarse slice of small categories with the corresponding homotopy theory of locally constant space-valued diagrams. The example above is the discrete-fibre part of that identification. It is essential that the coefficient diagram retain its transport maps; a diagram in the ordinary homotopy category specified only by the fibre objects does not supply this information.
\end{dcomment}

\begin{dcomment}{exposition}{p179}{Why a replacement may solve only one localization problem}
Grothendieck's replacements have three different tasks. The functor $\Phi$ makes the outgoing comma categories into fibres of a split fibration; $\Psi$ performs the incoming version; $C$ aims to make both directions harmless at once. Their usefulness depends on the class in which their comparison maps lie. The distinction between $\HOT_W^{\mathrm{lc}}$ and $\HOT_W^{\mathrm{lc}!}$ is already visible in those comparison maps.

We can isolate the categorical argument behind the repeated claims of equivalence. Let $\mathcal A\subseteq\mathcal B$ be a full subcategory, equipped with the restriction of a class $V$ in $\mathcal B$. Suppose there is a functor $R:\mathcal B\to\mathcal A$ and a natural transformation $\eta:\id_{\mathcal B}\to jR$, where $j$ is the inclusion, such that every $\eta_B$ lies in $V$. If $V$ has two-out-of-three, naturality shows that $R$ preserves $V$: for $u:B\to B'$, the equation
\[
 jR(u)\,\eta_B=\eta_{B'}\,u
\]
forces $jR(u)\in V$ whenever $u\in V$. After localization, $\eta$ gives an isomorphism $\id\simeq jR$. Its restriction to $\mathcal A$ gives $\id\simeq Rj$, so $j$ and $R$ induce quasi-inverse equivalences. This proof uses a natural comparison; objectwise existence of representatives alone would not control morphisms in the localization.

Apply the argument first to $\Phi$. On all of $\Cat/S$, the map $X\to\Phi(X)$ is in $W_S^{\mathrm d}$. It therefore proves that smooth representatives describe $\HOT_W(S)$. The split-fibration presentation uses in addition the cleavage retraction from the comma replacement of an already split fibration; that retraction respects cartesian transport. On $\Fib_W S$, the comparison is in $W_S^{\mathrm f}$ and the target belongs to the subcategory of split fibrations with weakly invertible transport. It therefore proves the corresponding statement about $\HOT_W^{\mathrm{lc}}(S)$. Dual statements hold for $\Psi$.

For $C$, the information available without W(4b) is instead
\[
 C(X)\in\Parf_\omega S,\qquad i_X:X\longrightarrow C(X)\in W_S.
\]
The replacement lemma applies perfectly well to \emph{coarse} localization. It identifies $(\Cat/S)W_S^{-1}$ with the coarse localization of $\Parf_\omega S$, and hence with the coarse localization of any intermediate full subcategory to which the same replacement applies. It says nothing comparable about the fine localization unless the comparison $i_X$ belongs to $W_S^{\mathrm f}$ when $X$ is already a weak fibration. That is the missing input, rather than a defect in the localization argument.

This is why $\lambda$ is a genuine additional functor:
\[
 \lambda:(\Fib_W S)(W_S^{\mathrm f})^{-1}
 \longrightarrow(\Fib_W S)(W_S)^{-1}.
\]
It is induced by enlarging the class of inverted maps. Calling both sides locally constant does not make $\lambda$ an equivalence. Its possible failure is different from the apparent multiplicity of intermediate \emph{fine} presentations. The manuscript's ``infinite series'' concerns a comparison that the present argument cannot establish; the alternating construction on pp.~192--195 subsequently resolves it.

In fact, under the revised W(1,2,3), \emph{every} full subcategory $\mathcal F$ with
\[
 \Parf_\omega S\subseteq\mathcal F\subseteq\Fib_W S
\]
has the same fine localization as $\Fib_W S$. To see the later answer, combine the automatic filtered-colimit closure of a fundamental localizer with the minimum theorem \citep[Corollary~4.2.22 and Theorem~6.1.18]{Cisinski2006}. They give W(6~bis) and identify the minimum obtained by adding this axiom to the definition of $W_\omega$: $W''_\omega=W_\omega=W_\infty$. The alternating replacement $R=\Phi_\infty$ therefore takes values in $\Parf_\omega S$, and its unit $X\to R(X)$ is fibrewise in $W$ whenever $X$ is a weak $W$-fibration. Applying the replacement lemma above, and then restricting $R$ to $\mathcal F$, gives
\[
 \mathcal F(W_S^{\mathrm f})^{-1}
 \xrightarrow{\sim}(\Fib_W S)(W_S^{\mathrm f})^{-1}.
\]
The proof of these properties of $R$ is supplied at p.~194, ``Why the alternating colimit repairs both directions''. Thus the proposed infinitely many inequivalent fine presentations do not survive the later results. This conclusion does not enlarge the class of maps inverted in the fine localization and does not make $\lambda$ an equivalence; the rational example below still distinguishes fine from coarse.

Under W(4b), both $X$ and $C(X)$ are weak fibrations whenever $X$ is, and $i_X\in W_S$ becomes $i_X\in W_S^{\mathrm f}$. The same replacement lemma then works with $C$ itself for fine localization. The additional effect of W(4b) is to identify the fine and coarse localizers on weak fibrations, reducing the four-category diagram to three. The alternating replacement achieves the equivalence of fine presentations without this extra identification. Restricting W(4b) to maps between already perfect objects does not supply the step for $i_X$ directly: its source is the object that has not yet been made perfect. The later indirect implication between the axioms is addressed at p.~184.

The unexplained W(4a$'$) at p.~180 can be identified from Chapter~XII, p.~53 \citep{GrothendieckDerivateursXII}. It says that pulling a $W$-equivalence back along a $\Parf_\omega$-morphism again gives a $W$-equivalence. W(4b), by contrast, tests the fibres of a map between two weak fibrations over one fixed base. The cartesian square displayed on pp.~179--180 explains the first route to $C(f)_s\in W$: its right and left sides are perfect, and its bottom side is $f$. Naturality of $i_X$ followed by W(4b) explains the second route. These are different assertions, even when a later theorem proves them equivalent under specified assumptions. The prime attached to the letter $W$ in W$'$(4a) of the earlier footnote should also be retained; it is not the same printed notation as W(4a$'$).
\end{dcomment}

\begin{dcomment}{exposition}{p182}{From replacement to reflection}
Grothendieck wants the locally constant objects to form a subcategory to which arbitrary relative homotopy types can be reflected. The equality of the two upper-floor localizers under W(4b) supplies their common category. The assertion $\rho\dashv i$ supplies more: it says what maps into that category mean.

It is helpful to spell out the universal property. Put
\[
 \mathcal D=(\Cat/S)(W_S^{\mathrm d})^{-1},\qquad
 \mathcal E=(\Cat/S)(W_S)^{-1}.
\]
The functor $\rho:\mathcal D\to\mathcal E$ comes from $W_S^{\mathrm d}\subseteq W_S$. The functor $i:\mathcal E\to\mathcal D$ is induced by $C$, since W(4b) turns $C(u)$ into a fine equivalence for every total equivalence $u$. The proposed unit at $X$ is the class of $i_X:X\to C(X)$. The adjunction asserts precisely that composing with this unit identifies maps from $X$ into any object in the image of $C$ with maps from its reflection. In the notation of the text, this is the second isomorphism in the boxed Hom formula; the first follows by the dual construction.

The formal verification uses the induced endofunctor $T=i\rho$ and its natural comparison $\eta:\id\to T$. W(4b) makes both $T\eta_X$ and $\eta_{TX}$ invertible: they are represented by $C(i_X)$ and $i_{C(X)}$, respectively. Naturality then gives the reflection property and forces the two iterates to agree. No independent coherence hypothesis is needed. The complete argument is given in the development at p.~197, ``Completing the reflection without a path calculation'', where it also resolves Grothendieck's later question about $\gamma i$ and $i\gamma$.

Once this universal property is available, the corollary for a weak fibration $Y$ is immediate in a precise sense. W(4b) gives $i_Y:Y\to C(Y)\in W_S^{\mathrm{univ}}$, so it is inverted both in $\mathcal D$ and in the analogous incoming-comma localization. Hence
\[
 \Hom_{W_S^{\mathrm d}}(X,Y)
 \xrightarrow{\sim}\Hom_{W_S^{\mathrm d}}(X,CY)
 \xrightarrow{\sim}\Hom_{W_S}(X,Y),
\]
and the same proof applies to $W_S^{\mathrm g}$. The bijections are induced by the localization functors themselves. They therefore identify the actual comparison maps in the corollary, rather than giving unspecified bijections between sets of the same size. Notice that only the target is required to be a weak fibration; the source can be any category over $S$. This asymmetry is exactly what makes a reflective subcategory useful for constructing maps from arbitrary coefficients.
\end{dcomment}

\begin{dcomment}{examples}{p184}{A rational equivalence with a nontrivial connected fibre}
Grothendieck proposes testing the distinction between the two locally constant categories with homological coefficients. The following calculation carries out that test over $\mathbb Q$. It shows that the difference is not merely a matter of connected components: a total equivalence can lose a positive-dimensional class in every connected fibre.

Let $W_{H\mathbb Q}$ be the class of functors inducing isomorphisms on the rational homology of nerves. This is an accessible fundamental localizer; the general construction from a homology theory is \citet[Theorem~9.4.2 and Corollary~9.4.3]{Cisinski2006}. Set
\[
 \Gamma=\langle a,t\mid tat^{-1}=a^{-1}\rangle
       =\mathbb Z\rtimes_{-1}\mathbb Z,
 \qquad S=B\mathbb Z,
\]
and let $p:B\Gamma\to S$ be induced by $a\mapsto0$, $t\mapsto1$. Write $X=(B\Gamma,p)$ and $Y=(S,\id_S)$ as objects of $\Cat/S$; the same functor gives $u:X\to Y$ over $S$.

We first check the relative hypotheses. Every integer $n$ has the lift $t^n$ in $\Gamma$. Since both categories are groupoids, any lift is cartesian and cocartesian: the universal factorization of an arrow is obtained by composing with its inverse. Thus $p$ is both a Grothendieck fibration and a Grothendieck opfibration. Its strict fibre is the one-object groupoid $B\langle a\rangle=B\mathbb Z$. Transport by the generator acts on that fibre by $a\mapsto a^{-1}$, an automorphism. The identity functor $Y\to S$ is also a bifibration, with terminal fibre. These objects belong even to $\Parf_\omega S$, since bifibrations are smooth and proper for the categorical asphericity structure and hence for each localizer under consideration.

We now compute the total rational homology without a spectral-sequence convergence issue. Let $a$ and $t$ act on $\mathbb R^2$ by
\[
 a(x,y)=(x+1,y),\qquad t(x,y)=(-x,y+1).
\]
They satisfy the relation above. A nonidentity element $a^m t^n$ has no fixed point: if $n\ne0$ it changes the second coordinate, and if $n=0$, $m\ne0$, it translates the first coordinate. The action is properly discontinuous, and its quotient is the usual square model of the Klein bottle. Since $\mathbb R^2$ is contractible, this quotient is a $K(\Gamma,1)$ and computes the homology of $B\Gamma$.

Its cellular structure has one vertex, the two oriented edges $a,t$, and one two-cell attached by $tat^{-1}a$. The cellular boundary of the two-cell records exponent sums, so its integral boundary is $2a$. Over $\mathbb Q$ the complex is therefore
\[
 0\longrightarrow\mathbb Q
 \xrightarrow{\ d_2\ }\mathbb Q a\oplus\mathbb Q t
 \xrightarrow{\ 0\ }\mathbb Q\longrightarrow0,
 \qquad d_2(1)=2a.
\]
The map to $B\mathbb Z\simeq S^1$ is induced by $p$, or geometrically by the second coordinate modulo~1. On cellular chains it sends $a$ to zero and $t$ to the circle generator. Its effect on $H_0$ and $H_1$ is an isomorphism, while both source and target have zero rational homology in higher degrees. Thus $u\in W_{H\mathbb Q,S}$.

The strict fibre of $u$ is $B\mathbb Z\to e$. It sends $H_1(B\mathbb Z;\mathbb Q)=\mathbb Q$ to zero, and is not in $W_{H\mathbb Q}$. Consequently
\[
 u\in W_S\setminus W_S^{\mathrm f}
 \quad\text{with }X,Y\in\Parf_\omega S.
\]
This disproves W(4b) and even the original W(4b$'$) for this localizer, exactly as Grothendieck's parenthesis anticipates. It also proves that $\lambda$ is not an equivalence in this example. Indeed, fibre homology defines a functor out of the fine localization, because it inverts all fibrewise rational equivalences. It detects that the class of $u$ is not invertible there. Coarse localization inverts $u$ by construction. Thus $\lambda$ fails to reflect isomorphisms.

The calculation explains the loss of information. The fibre has a one-dimensional rational homology group on which transport acts by $-1$. Both the invariant and coinvariant contributions of this action vanish because $1-(-1)=2$ is invertible. The total category consequently has the same rational homology as the base, although the fibre is not rationally aspherical. With coefficients of characteristic~2 the displayed differential vanishes and the total map is no longer a homology equivalence. The coefficient system changes the conclusion.

Cisinski's Proposition~9.4.6 establishes a related failure of properness for rational homology localization; that terminology concerns the associated model structures and should not be substituted without explanation for the manuscript's axiom W(4b). The elementary calculation above directly verifies the precise axiom used here. A later example at p.~198 shows the same obstruction already with finite discrete fibres and makes its effect on ordinary base change especially transparent.
\end{dcomment}

\begin{dcomment}{status}{p184}{The weakened fibre axiom after the later theorems}
Grothendieck distinguishes what his current replacement argument can prove from what a weaker axiom might eventually imply. The distinction is historical, and the footnote referring forward to Chapter~XVI already changes its status. We can now complete that update.

For a fundamental localizer, stability under filtered colimits is a theorem, with no accessibility assumption: \citet[Corollary~4.2.22]{Cisinski2006}. The extra limit hypothesis in the manuscript's forward reference is therefore available under the revised W(1,2,3). Applying Chapter~XVI, Theorem~4.4.13(2), manuscript pp.~58--60 \citep{GrothendieckDerivateursXVI}, with $W_0=W_\omega$, identifies the relevant conditions~B and~D: the original W(4b$'$), restricted to $\Parf_\omega$, is equivalent to W(4b) in this setting. The difficulty on pp.~179--184 identifies a missing argument in the proposed proof, not a surviving counterexample to that implication. Our rational example violates both axioms and is consistent with this equivalence.

There is a further restriction on the apparent range of the fibre axiom. Given a $W$-locally constant functor $u:A\to B$ in the sense that all transports $A/b\to A/b'$ lie in $W$, its comma-category replacement is a split opfibration over $B$ with fibres $A/b$. It is a weak $W$-fibration, and its total category is $W$-equivalent to $A$. If $u\in W$, W(4b), applied to the map from this replacement to $B\to B$, implies $A/b\to e\in W$ for every $b$. Thus a $W$-locally constant $W$-equivalence is $W$-aspherical.

The characterization in \citet[Theorem~6.4.22 and Remark~6.4.23]{Cisinski2006} says that the only nontrivial fundamental localizer with this property is $W_\infty$. Here ``nontrivial'' excludes the class of all functors. Conversely, ordinary weak homotopy equivalences satisfy the fibre axiom by the homotopy-fibre description of locally constant functors. It follows in particular that $W'_\omega=W_\omega=W_\infty$. These identifications settle the minimum and fibre-detection questions raised in this passage. They do not remove the interest of the four-category construction for general homological localizers, where W(4b) fails and the comparison $\lambda$ can genuinely lose information.
\end{dcomment}

% Commentary for manuscript pp.185--204.
\begin{dcomment}{exposition}{p185}{Changing the category on which an axiom is tested}
Grothendieck is trying to make the passage from a total category to its family of fibres as economical as possible. The replacements already constructed let him change the objects on which a localization is calculated. They do not, without further argument, let him change the arrows which that localization inverts. This is the distinction behind the long discussion of W(4b): a supply of perfect representatives and a theorem detecting equivalences on their fibres are different resources.

Here is the precise formal point in the paragraph on strong saturations. Suppose that $\mathcal A\subseteq\mathcal F$ is one of the replacement subcategories and that its inclusion induces equivalences both for the fine class $V=W_S^{\mathrm f}$ and for the coarse class $U=W_S$. Then the square
\[
\begin{tikzcd}
\mathcal A[V^{-1}]\arrow[r,"\sim"]\arrow[d]&\mathcal F[V^{-1}]\arrow[d]\\
\mathcal A[U^{-1}]\arrow[r,"\sim"]&\mathcal F[U^{-1}]
\end{tikzcd}
\]
commutes. If $U=V$ on $\mathcal A$, its left vertical arrow is an equivalence, so the right vertical arrow is one too. An arrow of $\mathcal F$ therefore becomes invertible in one localization exactly when it becomes invertible in the other. This is equality of the \emph{strongly saturated} classes in $\mathcal F$. Literal equality of $U$ and $V$ on $\mathcal F$ needs a separate saturation or detection hypothesis. The overlines in the edition are consequently essential, rather than an embellishment of the formula.

The proposed cartesian-functor condition is well suited to this argument because a split Grothendieck fibration is an actual diagram of categories. Its being locally constant means that every transport functor belongs to $W$. A cartesian map between two such fibrations is then a transformation of the diagrams, and its fibre maps are the components. W(4b) asks whether a weak equivalence of the two total categories forces all those components to be weak equivalences. The rational-homology example accompanying ms.~184 shows why the answer depends on the chosen notion of equivalence.

The later annotations must be read with the later axiom dictionary. Chapter XVI, Lemma~4.4.9, Corollary~4.4.10 and Theorem~4.4.13 organize the transfer between testing on split fibrations, weak fibrations and perfect maps, with the stated limit hypotheses \citep{GrothendieckDerivateursXVI}. They explain why an apparently weaker test can recover the original condition. The present discussion is not evidence that these conditions are inequivalent: it records the search for a proof of their relationship. Conversely, an equivalence obtained under countable-limit hypotheses does not remove those hypotheses from the argument written here.
\end{dcomment}

\begin{dcomment}{status}{p186}{Perfect maps and strict bifibrations}
Grothendieck wants a representative on which transport can be calculated in either direction. A perfect map provides the homotopical version of that symmetry. A bifibration provides actual cartesian and cocartesian lifts, with their universal properties. For an arrow $u:s\to t$, a chosen bifibration has an adjunction $u_!\dashv u^*$ between its fibres. Such an adjunction induces inverse homotopy equivalences on nerves. Thus strict bifibrations are excellent locally constant representatives; the question is whether there are enough of them.

A useful case can be settled without a new replacement construction. Let $F:S^{\mathrm{op}}\to\Cat$ send every arrow to $W$, and suppose that $S$ has a terminal object $e$. The arrows $s\to e$ give a natural transformation
\[
\operatorname{const}(F(e))\longrightarrow F,
\qquad F(e)\longrightarrow F(s),
\]
all of whose components lie in $W$. Its Grothendieck construction is a fibrewise weak equivalence from $S\times F(e)$ to the fibration classified by $F$. The projection $S\times F(e)\to S$ is a bifibration: both kinds of lifts keep the second coordinate fixed. If $S$ has an initial object $0$, the same argument uses $F\to\operatorname{const}(F(0))$. Applying the previously available split-fibration replacement first proves essential surjectivity of the strict-bifibration presentation on these bases. This proves the assertion about objects; it does not by itself establish full faithfulness after localization.

The perfect-representative question has a more general answer in the manuscript itself: under W(6 bis), the alternating construction below gives a fibrewise weak equivalence into a perfect object. In the usual weak-homotopy theory, later results give a functorial factorization through a smooth and proper map, and describe the resulting category of fibrant objects \citep[Proposition~4.4]{Cisinski2009LocallyConstant}. A modern space-valued version identifies locally constant right fibrations with Kan fibrations \citep[Theorem~4.6.10, author's July~2026 revision]{Cisinski2019HigherCategories}. Both statements realize the desired homotopical symmetry.

They must nevertheless be distinguished from the literal question about $\Bifib S$ in this paragraph. A smooth and proper functor need not be a Grothendieck bifibration, and a Kan fibration is a simplicial object, not a functor of ordinary categories with chosen universal lifts. The recent account of categorical fibre bundles starts with transports that are categorical isomorphisms; its bifibration theorem does not supply replacement of an arbitrary diagram with $W$-equivalence transports \citep[\S3 and Theorem~4.5]{CarcaciaCampos2026}. In the checked works of Maltsiniotis and Cisinski we could not find a theorem identifying the precise strict-bifibration localization above for every $W$ under only the axioms retained here. We therefore leave that assertion unsettled, while separating it from the perfect-replacement theorem which the following pages establish. This is a boundary of the present verification, not a declaration that the question is an open problem in every usual homotopy theory.
\end{dcomment}

\begin{dcomment}{examples}{p190}{A fork which defeats a one-sided replacement}
Grothendieck's example isolates the difficulty of producing transport in both directions. Over an interval, a terminal fibre can govern the total homotopy type; over a fork, the two outgoing branches can carry different information. We spell out his construction because it proves the failure directly, without appealing to the suggested amalgamated sum.

Let $S$ be the poset $0<a$, $0<b$, with $a$ and $b$ incomparable. Choose a category $T$ and an object $t\in T$, and let $X\to S$ be the split opfibration classified by
\[
T\longleftarrow e\xrightarrow{\mathrm{id}}e,
\]
where the left functor selects $t$. Besides the objects of $T$ over $a$, the total category $X$ has an object $z$ over $0$ and an object $v$ over $b$. There is one arrow $z\to v$; the arrows $z\to y$, for $y\in T$, are the morphisms $t\to y$ in $T$; all other nonidentity arrows are those already in $T$. Composition of an arrow out of $z$ with an arrow in $T$ is composition in $T$. There are no arrows between $v$ and the objects of $T$.

Write $X''$ for the full subcategory on $z$ and $T$. A retraction $r:X\to X''$ fixes $X''$ and sends $v$ to $z$, and sends $z\to v$ to $\mathrm{id}_z$. If $j:X''\hookrightarrow X$, then $rj=\mathrm{id}$ and there is a natural transformation $jr\to\mathrm{id}_X$, equal to the identity on $X''$ and to $z\to v$ at $v$. Naturality has only one nonidentity arrow outside $X''$ to check, namely $z\to v$, for which both composites are that arrow. Hence $j$ is a homotopy equivalence of categories.

There is also a retraction $q:X''\to T$, sending $z$ to $t$ and the arrow $z\to y$ labelled by $h:t\to y$ to $h$. For the inclusion $k:T\hookrightarrow X''$ we have $qk=\mathrm{id}_T$. The arrow $z\to t$ labelled by $\mathrm{id}_t$, together with the identity arrows on $T$, defines a natural transformation $\mathrm{id}_{X''}\to kq$. For an arrow $z\to y$ labelled by $h$, its naturality square is precisely the equation $h\mathrm{id}_t=h$. Thus $k$ too is a homotopy equivalence. The composite $T\to X$ lies in every fundamental localizer. This verifies the proposed contraction, including its functoriality. The edition reverses $X$ and $X''$ in the sentence about adjoining $v$; the retraction that follows makes the intended direction unambiguous.

The fibres of the split fibration $\Phi_S(X)$ are
\[
\Phi_S(X)_0=0\backslash X\cong X,
\qquad \Phi_S(X)_a\cong T,
\qquad \Phi_S(X)_b\cong e.
\]
Its transport along $0\to b$ is the inclusion $e\to X$ selecting $v$. If $T$ is not $W$-aspherical, this map is not in $W$: otherwise $X\to e$ and then $T\to e$ would be in $W$ by two-out-of-three. Thus $X\to S$ is a Grothendieck opfibration, hence $W$-proper, while $\Phi_S(X)$ is not even a weak $W$-fibration. For $W=W_\infty$, one may take $T$ to be the discrete two-object category and select either object. Its two connected components make the obstruction visible without higher homotopy groups.

The same fork separates perfectness from local constancy even when all original fibres are points. Take $X=S$ with its identity map. Its replacement is the arrow category of $S$ projected by source. The fibres at $0,a,b$ are $S,e,e$, and both transport maps are weak equivalences because $S$ has an initial object. Yet the object corresponding to $a\in S=\Phi_S(S)_0$ has no lift towards the fibre over $b$: such a lift would require an arrow $a\to b$. The relevant lifting category is empty, so the projection is not $W$-proper for any nondegenerate $W$. One application of $\Phi$ can therefore destroy properness even for the identity of the base.

There is also a precise affirmative version of the aside about pushouts. For a span of full inclusions $A\hookrightarrow B$, $A\hookrightarrow C$ which are both sieves, a chain of composable arrows in the amalgamated category cannot pass from $B\setminus A$ through $A$ into $C\setminus A$: the first part is forbidden by the sieve condition. The same argument, with the last part forbidden, applies when both inclusions are cosieves. All chains therefore lie wholly in one side, and
\[
N(B\amalg_A C)=NB\amalg_{NA}NC.
\]
The two maps out of $NA$ are monomorphisms of simplicial sets, so this pushout computes their homotopy pushout. Thomason's homotopy-colimit theorem identifies the nerve of the Grothendieck construction of the span with that same homotopy pushout \citep[Theorem~1.2]{Thomason1979}. These are the explicit sufficient hypotheses needed for Grothendieck's aside. His particular contraction above is stronger in a different direction: it works for an arbitrary choice of $t\in T$, without requiring the inclusion of that object to be a sieve or a cosieve.
\end{dcomment}

\begin{dcomment}{development}{p192:Mais chemin faisant est apparue une idée bizarre.}{The support obstruction behind the fork}
Grothendieck's stronger assertion on the fork suggests that the obstruction comes from the base rather than from a badly chosen total category. The following extension makes that point precise. It is a direct elaboration of his lifting argument; we have not found this support criterion stated in the checked treatments of smooth and proper functors, and claim no priority for it.

\textbf{Proposition (incompatible upper cones).}\ours\
Let $P$ be a poset, let $s\leq a,b$, and suppose that no element of $P$ is greater than or equal to both $a$ and $b$. Let $p:X\to P$ be a functor, and let $W$ be a fundamental localizer for which $\varnothing\to e$ is not in $W$. If $\Phi_P(X)\to P$ is $W$-proper, then $s\backslash X$ is empty. In particular, if $s$ is initial in $P$, then $X$ is empty.

To prove the proposition we use the arrow-lifting criterion for properness. For a Grothendieck fibration classified by a contravariant diagram $F$, properness requires, for $u:s\to t$ and $z\in F(s)$, the category $z\backslash u^*$ to be $W$-aspherical. Its objects are pairs $(y,z\to u^*y)$, with $y\in F(t)$. This is the proper dual of the usual comma-category criterion for smoothness \citep[Theorem~4.13 and Example~4.17]{Maltsiniotis2005Asphericite}. In particular this category must be nonempty.

Here $F(t)=t\backslash X$. Suppose, for a contradiction, that $F(s)$ has an object. Applying the criterion to $s\leq a$ produces an object $(x,a\leq p(x))$ of $F(a)$. Regard it also as an object $z$ of $F(s)$ by restriction. Apply the same criterion to $s\leq b$ and this particular $z$. It produces an object $(y,b\leq p(y))$ of $F(b)$ and an arrow $z\to (s\leq p(y),y)$ in $F(s)$. The latter contains an arrow $x\to y$ in $X$. Consequently
\[
a\leq p(x)\leq p(y),\qquad b\leq p(y),
\]
contradicting the assumed disjointness of the two upper cones. Hence $F(s)$ is empty. If $s$ is initial, $s\backslash X$ is canonically $X$, which proves the last assertion.

For the original fork, the two upper cones are $\{a\}$ and $\{b\}$. The proof applies to every $X\to S$, including categories with objects only over $0$. This fills the last step of the displayed assertion on ms.~192: after the two end fibres have been eliminated, coasphericity of the empty end-fibre inclusion also eliminates the remaining total category. No locally constant hypothesis on $X$ is needed.

The nondegeneracy assumption has a concrete role. An empty lifting category cannot be aspherical. If $\varnothing\to e$ belonged to a product-stable fundamental localizer, then $\varnothing\to A$ would belong to it for every $A$, and two-out-of-three would put every functor in $W$. The exceptional case excluded by Grothendieck is therefore exactly the case in which this obstruction disappears.
\end{dcomment}

\begin{dcomment}{exposition}{p194}{Why the alternating colimit repairs both directions}
Grothendieck now replaces the search for a single miraculous functor by a process in which every finite stage will eventually be improved in the other direction. There need be no finite stage which is both smooth and proper. What matters is that the smooth stages and the proper stages each form a cofinal subsequence. The extra limit axiom lets their two properties survive in the same colimit.

We first justify the closure assertion used here, keeping the scope of W(6 bis) explicit. For a sequential diagram of functors $A_n\to B_n$, a chosen object and arrow in its colimit are represented at some finite stage. The category of lifts of that arrow is the filtered colimit, over the remaining stages, of the corresponding lifting categories. Objects, arrows, commutativity requirements and equalities of arrows involve only finitely many data, so every such datum and every equality is witnessed at a stage. If all $A_n\to B_n$ are $W$-proper, these lifting categories are $W$-aspherical. Apply W(6 bis) to their maps to $e$: their colimit is $W$-aspherical as well. The outgoing-arrow lifting criterion proves that the colimit functor is proper. Reversing arrows gives the assertion for smooth functors. This proof uses sequential, or countably cofinal filtered, diagrams. The unqualified phrase about filtered colimits in the edition abbreviates this scope. A fully general filtered-colimit theorem needs the corresponding closure of $W$.

Put $X_0=X$ and let $X_1,X_2,\ldots$ be the alternating sequence displayed in the text. Every odd stage is of the form $\Phi(Z)$ and is smooth over $S$; every positive even stage is of the form $\Psi(Z)$ and is proper. The odd indices and the positive even indices are both cofinal in $\mathbf N$. Consequently $R(X)=\varinjlim_nX_n$ is the colimit of a sequence of smooth maps to $S$, and also the colimit of a sequence of proper maps to $S$. The preceding closure argument makes $R(X)\to S$ perfect.

The assertion about the comparison map is equally concrete. Suppose first that $X\to S$ is a weak $W$-fibration. All $X_n$ remain weak $W$-fibrations and every transition is fibrewise in $W$. For each $s$, the maps $X_s\to(X_n)_s$ are therefore in $W$. Strict fibres commute with filtered colimits of categories: the equations specifying that an object or arrow lies over $s$ are finite equations, and hence are witnessed at a stage. Applying W(6 bis) to the transformation from the constant diagram $X_s$ to the fibre sequence gives
\[
X_s\longrightarrow\varinjlim_n(X_n)_s=R(X)_s\quad\text{in }W.
\]
For an arbitrary $X$, apply the same argument to the total categories. Each finite composite $X\to X_n$ is in $W$, so $X\to R(X)$ is a coarse weak equivalence. It need not be fibrewise; nothing in the argument asserts that it is.

These comparison maps also verify the localization claim. If $u:X\to Y$ is a fibrewise weak equivalence between weak fibrations, naturality gives
\[
R(u)\,i_X=i_Yu.
\]
Both $i_X$ and $i_Y$ are fibrewise in $W$, so $R(u)$ is fibrewise in $W$ by two-out-of-three. Thus $R$ descends to the fine localization. Its restriction to perfect objects has the same natural weak-equivalence unit, so the inclusion of perfect objects and the functor induced by $R$ are quasi-inverse equivalences. The argument works unchanged for any full intermediate subcategory containing the image of $R$. The coarse version uses the same equation and the coarse comparison maps. This proves the assertions on ms.~195 without silently replacing the fine class by the coarse one.

The later theory gives these constructions an independent setting. The infinite-path factorization in \citet[Sections~3.2.40--3.2.45]{Maltsiniotis2020} provides explicit smooth and proper replacements. Moreover, a modern fundamental localizer is stable under all filtered colimits \citep[Corollary~4.2.22]{Cisinski2006}; W(6 bis) is therefore automatic in that setting. The extra axiom has a real role in the proof written here, but is not an additional restriction in the completed modern theory. Together with the minimal-localizer comparison discussed earlier, this also removes the apparent distinction between the two minimal classes used for the perfect-representative question on ms.~186.
\end{dcomment}

\begin{dcomment}{development}{p197:Je ne le ferai que sous la pression des besoins.}{Completing the reflection without a path calculation}
Grothendieck has reached a precise remaining verification: applying the replacement to its unit and taking the unit of the replacement should give the same map after fine localization. His proposed geometric comparison can be replaced by a short categorical argument. We give the details because they complete both the present proposition and the Hom formula on ms.~182. The general setting is reflective localization, or equivalently idempotent monads \citep[Proposition~5.3.3]{Riehl2016}; the application below is our explicit completion of the omitted verification, not a claim of a new general localization theorem.

\textbf{Completion of the reflection argument.}\ours\
Under the comparison hypotheses of the proposition on ms.~196, both replacements specified there satisfy $i_{\gamma X}=\gamma(i_X)$ after colocal localization and give the asserted reflection. The formal tool is the following elementary lemma, a standard form of reflective localization.

\textbf{Lemma (a coaugmented functor).}
Let $T:\mathcal A\to\mathcal A$ have a natural transformation $\eta:\mathrm{id}\to T$. Suppose that $T\eta_X$ and $\eta_{TX}$ are isomorphisms for every $X$. Let $\mathcal A_{\mathrm{loc}}$ be the full subcategory of objects $Y$ for which $\eta_Y$ is an isomorphism. Then $T$, regarded as taking values in $\mathcal A_{\mathrm{loc}}$, is left adjoint to its inclusion, with unit $\eta$. In addition,
\[
T\eta_X=\eta_{TX}:TX\longrightarrow T^2X.
\]

Indeed, $TX$ is local by hypothesis. For local $Y$, consider precomposition
\[
\Hom_{\mathcal A}(TX,Y)\longrightarrow\Hom_{\mathcal A}(X,Y),
\qquad g\longmapsto g\eta_X.
\]
Given $f:X\to Y$, put $g=\eta_Y^{-1}T(f)$. Naturality of $\eta$ gives $T(f)\eta_X=\eta_Yf$, so $g\eta_X=f$. This proves surjectivity. If $g\eta_X=h\eta_X$, apply $T$ and cancel the isomorphism $T\eta_X$ to obtain $Tg=Th$. Naturality now gives
\[
\eta_Yg=Tg\,\eta_{TX}=Th\,\eta_{TX}=\eta_Yh.
\]
Since $\eta_Y$ is invertible, $g=h$. We have proved a natural bijection of Hom sets, which is the adjunction. Finally, naturality at $\eta_X$ says that $T\eta_X$ and $\eta_{TX}$ have the same composite with $\eta_X$. Their target $T^2X$ is local, so injectivity of the same Hom map proves their equality. This proves the lemma, including the coherence which initially appeared to require a separate calculation.

Apply it to
\[
\mathcal A=(\Cat/S)[(W_S^{\mathrm d})^{-1}],
\qquad T=\gamma,
\]
where $\gamma=C$ or $\Phi_\infty$ is now regarded as an endofunctor of the fine localization. The hypotheses of the proposition ensure that $\gamma$ sends coarse weak equivalences to arrows invertible in this fine localization. Since $i_X$ is a coarse weak equivalence, $\gamma(i_X)$ is invertible. The other map $i_{\gamma X}$ is a coarse equivalence between weak $W$-fibrations. Under the proposition's comparison hypotheses, $\lambda$ is an equivalence, so this map is already invertible in the fine localization of weak fibrations and hence in the colocal localization displayed above. For $\gamma=\Phi_\infty$, the previously established fibrewise-unit property gives this conclusion directly. For $\gamma=C$, the detection or comparison hypothesis is essential; weak fibrancy alone would not justify it. Thus both hypotheses of the lemma hold, and its conclusion is exactly
\[
i_{\gamma X}=\gamma(i_X).
\]
There is no cancellation of $i_X$ itself in the fine localization: that arrow is generally not invertible there. The cancellation occurs only after applying $\gamma$, where its invertibility is known. This distinction is the whole reason that a direct appeal to naturality alone would have been insufficient.

The local objects are precisely the essential image of the fine locally constant category. Every such representative has invertible unit, while an object with invertible unit is isomorphic to its replacement, which lies in that image. Moreover, $T$ inverts coarse equivalences and its unit is a coarse equivalence, so it has the same localization universal property as $\rho$. The reflector therefore identifies with $\rho$, and its fully faithful right adjoint with $i$. The dual argument gives $\sigma\dashv j$. In particular, for a local object $L$, a map $X\to L$ factors uniquely through $X\to\gamma X$; this is the Hom description asserted earlier. The verification applies to the alternating replacement just as it does to the path replacement, and closes the expressly postponed step.
\end{dcomment}

\begin{dcomment}{examples}{p198}{A coarse equivalence lost on taking one fibre}
Grothendieck's warning about coarse restriction concerns the most elementary base change: passage to a single object of the base. A finite group supplies a compact example, complementary to the connected-fibre example at ms.~184.

Let $G$ be a finite group with more than one element, and let $BG$ denote its one-object category. Define $E$ to have object set $G$ and exactly one arrow between every ordered pair of objects. Send its unique arrow $g\to h$ to $hg^{-1}\in G$. This defines a functor $p:E\to BG$, since $(kh^{-1})(hg^{-1})=kg^{-1}$. Every arrow of $BG$ lifts uniquely once either its source or target in $E$ is prescribed, so $p$ is both a discrete fibration and a discrete opfibration. Its strict fibre over the unique base object is the discrete category $G$, and its transports are permutations.

Take $W$ to be the functors inducing isomorphisms on rational homology of nerves. The category $E$ is contractible: every object is both initial and terminal. The rational homology of $BG$ is $\mathbf Q$ in degree zero and zero in positive degrees. One algebraic verification uses the averaging idempotent
\[
e_G=|G|^{-1}\sum_{g\in G}g\in\mathbf Q[G].
\]
The trivial $\mathbf Q[G]$-module is the direct summand $\mathbf Q[G]e_G$ of the regular module, hence is projective. The higher groups $\operatorname{Tor}^{\mathbf Q[G]}_n(\mathbf Q,\mathbf Q)$ vanish, giving the asserted group homology. Thus $p:E\to BG$ is in $W$.

Regard this as a morphism over $BG$ from $E\to BG$ to $BG\xrightarrow{\mathrm{id}}BG$. Both objects are perfect and weakly $W$-fibred; all their transports are isomorphisms. Pull back along the point $e\to BG$. The resulting map is
\[
G_{\mathrm{disc}}\longrightarrow e.
\]
On degree-zero rational homology it is $\mathbf Q^{|G|}\to\mathbf Q$, sending every basis vector to $1$, and is not an isomorphism. Hence raw pullback does not descend through the coarse class even on these excellent representatives. The fine class poses no problem: the original map was never fibrewise in $W$.

This also explains why a functor on the coarse localization cannot be manufactured merely by observing that all its objects have fine replacements. The obstruction is a specific arrow which has already been inverted. Any proposed pullback agreeing with strict fibre on these representatives would have to send that invertible arrow to a noninvertible one. Under W(4b) this example is excluded at the outset; without it, the distinction between the two localizations records precisely this lost fibre information.
\end{dcomment}

\begin{dcomment}{exposition}{p201}{An endofunctor identity and a representative formula}
Grothendieck has already obtained the right endofunctor identity before settling the formula for an arbitrary pair of raw representatives. For fixed $X$, the construction $\mathrm{mult}_{X/S}$ is defined by evaluating $X\times_S(-)$ on smooth representatives of its input. Therefore proving
\[
\mathrm{mult}_{X/S}(\eta)=\mathrm{hot}_{S,W}(X)\times_S\eta
\]
requires a calculation only on smooth representatives of $\eta$. It does not permit evaluating that functor on an arbitrary raw $Y$ by the formula $\mathrm{hot}_{S,W}(X\times_S Y)$.

This explains the different scopes of $(*)$ and $(**)$ on ms.~201. The argument with $X\to\Phi X$ establishes $(**)$ for a weak $W$-fibration $X$, because the second input may be chosen smooth. The further claim $(*)$ requires compatibility of multiplication by $X$ with replacement of the second raw input. Properness of $X$ is one sufficient route, but it is stronger than what this particular comparison needs. In the next block we prove that the locally constant transport available after replacing a weak $W$-fibration gives the missing compatibility directly. Thus the tentative negative expectation at the end of ms.~202 is not the final mathematical answer.
\end{dcomment}

\begin{dcomment}{development}{p202:Ça doit être faux sans hypothèse de}{The raw product with a weakly fibred factor}
Grothendieck seeks a condition under which an ordinary fibre product already represents the product of homotopy types. His attempted argument asks for a colocal equivalence after pulling back to the total category $X$. For the desired conclusion it is enough to obtain a colocal equivalence after viewing the resulting map over $S$. Locally constant transport supplies this weaker conclusion. The following proof completes his unrestricted formula $(*)$ on ms.~201 under the standing basic localizer hypotheses. We did not find this exact statement about raw representatives in the checked discussions of locally constant diagrams; the proof uses the replacement and localizer properties already established here, and carries no claim of priority.

\textbf{Proposition (a locally constant factor).}\ours\
Assume W(1,2,3), with the smooth and weak-fibration replacement properties used in these pages. If $X\to S$ is a weak $W$-fibration and $Y\to S$ is arbitrary, the canonical comparison
\[
\mathrm{hot}_{S,W}(X\times_S Y)
\longrightarrow
\mathrm{hot}_{S,W}(X)\times_S\mathrm{hot}_{S,W}(Y)
\]
is an isomorphism. No W(4b), W-properness of $X$, or smoothness of $Y$ is required.

We first prove the key assertion for a split Grothendieck fibration $E\to S$ classified by $F:S^{\mathrm{op}}\to\Cat$, all of whose transport functors are in $W$. Namely,
\[
v:Y\to Z\text{ in }W_S^{\mathrm d}
\quad\Longrightarrow\quad
E\times_S v\text{ in }W_S^{\mathrm d},
\]
where the latter map is regarded as a map \emph{over $S$}. Fix $s\in S$ and put $I_s(Y)=s\backslash Y$. An object is $(y,\alpha:s\to p(y))$, and a morphism to $(y',\alpha')$ is $h:y\to y'$ satisfying $p(h)\alpha=\alpha'$. The category $s\backslash(E\times_S Y)$ is the Grothendieck fibration over $I_s(Y)$ with fibre $F(p(y))$ at $(y,\alpha)$. Its transition along $h$ is $F(p(h))$.

The functors $F(\alpha):F(p(y))\to F(s)$ constitute a natural transformation to the constant diagram with value $F(s)$. Indeed,
\[
F(\alpha)F(p(h))=F(p(h)\alpha)=F(\alpha').
\]
They are all in $W$, since $F$ is locally constant. Passing to Grothendieck constructions gives a functor
\[
H_{s,Y}:s\backslash(E\times_S Y)
\longrightarrow F(s)\times(s\backslash Y)
\]
which is fibrewise in $W$ over $I_s(Y)$. Both source and target are Grothendieck fibrations over that category, hence smooth. The previously proved implication from fibrewise to global equivalence for smooth objects therefore shows that $H_{s,Y}$ belongs to $W$. Explicitly it sends $(x,y,\alpha)$ to $(F(\alpha)x,(y,\alpha))$; the displayed naturality equation defines its action on arrows.

This construction is natural in $Y$. For $v:Y\to Z$ over $S$, we obtain a commutative square
\[
\begin{tikzcd}[column sep=small]
s\backslash(E\times_S Y)\arrow[r]\arrow[d,"H_{s,Y}"']&
s\backslash(E\times_S Z)\arrow[d,"H_{s,Z}"]\\
F(s)\times(s\backslash Y)\arrow[r]&F(s)\times(s\backslash Z).
\end{tikzcd}
\]
The vertical maps are in $W$. If $v\in W_S^{\mathrm d}$, the lower map is in $W$, because $s\backslash v$ is in $W$ and $W$ is stable under product with a category. Two-out-of-three puts the upper map in $W$. This holds for every $s$, proving the key assertion.

Return to arbitrary weakly $W$-fibred $X$. Its replacement $E=\Phi_S(X)$ is a split fibration with locally constant transport. The map $i_X:X\to E$ is fibrewise in $W$. After base change to any $Y$, it remains fibrewise in $W$ between weak $W$-fibrations over $Y$; hence it is colocal over $Y$. Postcomposition $Y\to S$ preserves colocal equivalences, so
\[
X\times_S Y\longrightarrow E\times_S Y
\quad\text{belongs to }W_S^{\mathrm d}.
\]
Now $i_Y:Y\to\Phi_S(Y)$ belongs to $W_S^{\mathrm d}$. The key assertion applied to $E$ gives
\[
E\times_S Y\longrightarrow E\times_S\Phi_S(Y)
\quad\text{in }W_S^{\mathrm d}.
\]
The last object is the fibre product of two smooth representatives, so its class is the product of their classes. Composing these two equivalences proves the proposition, and the maps are precisely the canonical replacement comparison.

There is no contradiction with the fork example or the characterization of proper maps. The argument establishes membership in $W_S^{\mathrm d}$ after postcomposition to $S$. It does not establish membership in $W_E^{\mathrm d}$ when the same product map is regarded over $E$. Properness would supply that stronger compatibility. The distinction permits a complete answer to the product question while leaving the failure of properness entirely intact.
\end{dcomment}

\begin{dcomment}{later}{p204}{What the final fragment has established}
Grothendieck ends this part with a calculus whose two essential operations have concrete categorical representatives. Direct image is postcomposition; inverse image is pullback after a smooth replacement. Their adjunction and the comma-category formula for the fibres of direct image have been proved. The product computation above shows, in addition, that a weakly fibred factor suffices for the raw product formula which he had left doubtful.

The last displayed square is followed by an explicit notice that the continuation is missing. We do not reconstruct lost prose. The mathematical conclusion suggested by the earlier theorem is nevertheless determined: for a \emph{cartesian} square with the displayed notation, the mate
\[
f'_!q^*\longrightarrow g^*f_!
\]
is an isomorphism when $f$ is $W$-smooth or $g$ is $W$-proper. In the first case, a smooth representative over $X$ remains smooth after composing with $f$, so both sides are represented by the same pullback. In the second, proper base change computes restriction on arbitrary representatives, and the same calculation applies. These are the arguments already given on ms.~162--169, now under the revised localizer conventions. The missing conclusion belongs in commentary, not in the authoritative transcription.

The list of axioms on ms.~204 is deliberately shorter than a complete modern derivator axiom system. In particular, pointwise detection of isomorphisms has not been proved merely by writing the adjunction and pointwise formula. Its general later repair for fundamental localizers is explained at ms.~165. Model-category and derivator constructions supply the surrounding diagram calculus; the modern numbering and its distinction from Grothendieck's Der~3 conditions should be kept explicit \citep[Section~1]{Groth2013}. For a fixed base category, locally constant objects and their diagrams in auxiliary shapes form a full subderivator with both adjoints to its inclusion; weak homotopy equivalences of the base induce equivalences on the corresponding locally constant theories \citep[Theorems~3.12 and~1.3]{Cisinski2009LocallyConstant}. The fixed base and the varying auxiliary shape have different roles in this statement. These later theorems do not, by themselves, identify the coarse presentation with the fine one: the finite-group rational-homology example shows why that comparison still needs the relevant detection hypothesis.
\end{dcomment}
